% Imported from LaTeX_and_pdfs/AGAT_dump.tex; see provenance.json.
\part{Category Theory}


\section{Note on notation}

From the section on ``Terminology'' of the Preface of Barr and Wells (1998) \cite{agat:BW1998}:

\begin{quote}
	``In most scientific disciplines, notation and terminology are standardized, of- ten by an international nomenclature committee. (Would you recognize Einstein’s equation if it said $p = HU^2$?) We must warn the nonmathematician reader that such is not the case in mathematics. There is no standardization body and terminology and notation are individual and often idiosyncratic.''
\end{quote}

To try to bridge the difference choice of notation and through comparison, suggest the ``best'' notation that's easy to remember and easy to use, I'll present all the different types of notation that I come across as much as I can. My plan of attack is the following:
\begin{enumerate}
	\item I'll try to present different types of notation and reference the authors of the text when I can.
	\item I'll try to defer to the notation used in Wikipedia, first.
	\item I'll make a final decision of what notation works best (for me).
\end{enumerate}

\section{Category $\mathbf{A}$, (definition)}

\begin{definition}[Category $\mathbf{A}$]
\textbf{category} $\mathbf{A}$ is quadruple $\mathbf{A} = (\text{Obj}(\mathbf{A}), \text{Mor}\mathbf{A}, 1,\circ)$ 
\begin{equation}
\mathbf{A} = (\text{Obj}(\mathbf{A}), \text{Mor}\mathbf{A}, 1,\circ)
\end{equation}
s.t.
\begin{enumerate}
	\item $\text{Obj}{(\mathbf{A})}$ is a \emph{class}, whose elements, $A \in \text{Obj}(\mathbf{A})$, are called \emph{objects}
	\item $\text{Mor}\mathbf{A}$ is a \emph{class}.
	\begin{enumerate}
		\item From Ad\'{a}mek, Herrlich, and Strecker (2004) \cite{agat:AHS2004}, Kashiwara and Schapira (2006) \cite{agat:KaSch2006}, \\
		$\forall \, A, B \in \text{Obj}(\mathbf{A})$, $\exists \, \text{Hom}(A,B) \subseteq \text{Mor}(\mathbf{A})$. Therefore, 
		\begin{equation}
		\text{Mor}\mathbf{A} = \bigcup_{ A,B \in \text{Obj}(\mathbf{A}) } \text{Hom}(A,B)
		\end{equation}
		\item $\forall \, f \in \text{Hom}(A,B)$, $f: A \to B \in \text{Hom}(A,B)$ is a \emph{morphism}. Leinster (2014) \cite{agat:Lein2014} also calls them on pp. 10 \textbf{maps} or \textbf{arrow} from $A$ to $B$.
	\[
\begin{tikzpicture}
\matrix (m) [matrix of math nodes, row sep=3em, column sep=4em, minimum width=2em]
{
	A & B  \\
};
\path[->]
(m-1-1) edge node [above] {$f$} (m-1-2)
;
\end{tikzpicture}
\]
	\end{enumerate}
	\item $\forall \, A \in \text{Obj}(\mathbf{A})$, $\exists \, 1_A : A \to A$, i.e. $\exists \, \mathbf{1}_A \in \text{Hom}_{\mathbf{A}}(A,A) \equiv \text{Hom}(A,A)$, 
	\[
	\begin{tikzpicture}
	\matrix (m) [matrix of math nodes, row sep=3em, column sep=4em, minimum width=2em]
	{
		A & A  \\
	};
	\path[->]
	(m-1-1) edge node [above] {$1_A$} (m-1-2)
	;
	\end{tikzpicture} \text{ or } \begin{tikzpicture}
	\matrix (m) [matrix of math nodes, row sep=3em, column sep=4em, minimum width=2em]
	{
		A   \\
	};
	\path[->]
	(m-1-1) edge [loop left] node [above] {$\mathbf{1}_A$} (m-1-1)
	;
	\end{tikzpicture}
	\]
	\item \textbf{composition:} 
	$\forall \, A,B,C \in \text{Obj}\mathbf{A}$, define \textbf{composition} to be a map
	\begin{equation}
	\begin{aligned}
		\text{Hom}_{\mathbf{A}}(A, B) \times \text{Hom}_{\mathbf{A}}(B, C) & \to \text{Hom}_{\mathbf{A}}(A, C) \\
		 (f, g) & \mapsto g\circ f
 	\end{aligned}
\end{equation}, i.e.
	
	$\forall \, \begin{aligned} & \quad \\
	& f: A \to B \in \text{Hom}(A,B) \\
	& g:B\to C \in \text{Hom}(B,C) \end{aligned}$, i.e. $f,g \in \text{Mor}\mathbf{A}$, \\
	then $g\circ f : A \to C \in \text{Hom}(A,C)$, $g\circ f \in \text{Mor}\mathbf{A}$ i.e. 
	\[
	\begin{tikzpicture}
	\matrix (m) [matrix of math nodes, row sep=3em, column sep=4em, minimum width=2em]
	{
		A & B & C  \\
	};
	\path[->]
	(m-1-1) edge node [above] {$f$} (m-1-2)
	(m-1-2) edge node [above] {$g$} (m-1-3)
	(m-1-1) edge [bend right=30] node [below] {$g\circ f$} (m-1-3) 
	;
	\end{tikzpicture}
	\]
	s.t. 
	\begin{enumerate}
		\item \emph{associativity} $\forall \, \begin{aligned} & \quad \\
		& f: A \to B \\
		& g: B \to C \\
		& h: C \to D \end{aligned}$, $h\circ (g\circ f) = (h\circ g) \circ f $ i.e.
		
		\[
		\begin{tikzpicture}
		\matrix (m) [matrix of math nodes, row sep=3em, column sep=4em, minimum width=2em]
		{
			A & B & C  & D \\
		};
		\path[->]
		(m-1-1) edge node [above] {$f$} (m-1-2)
		(m-1-2) edge node [above] {$g$} (m-1-3)
		(m-1-3) edge node [above] {$h$} (m-1-4)
		(m-1-1) edge [bend right=30] node [below] {$g\circ f$} (m-1-3) 
		(m-1-2) edge [bend left=30] node [above] {$h\circ g$} (m-1-4) 
		;
		\path[-{>[scale=1.15]}, line width=0.5pt]
		(m-1-1) edge [bend left=45] node [above] {$(h \circ g)\circ f$} (m-1-4)
		edge [bend right=45] node [below] {$h \circ (g\circ f)$} (m-1-4);
		\end{tikzpicture}
		\]
		
		\item $\forall \, f:A \to B \in \text{Hom}(A,B)$, $1_B \circ f = f $ and $f\circ 1_A = f$ i.e.
		
		$\forall \, f \in \text{Hom}_{\mathbf{A}}(A,B)$,
		\[
		\begin{tikzpicture}
		\matrix (m) [matrix of math nodes, row sep=3em, column sep=4em, minimum width=2em]
		{
			A & B  \\
		};
		\path[->]
		(m-1-1) edge node [above] {$f$} (m-1-2)
		edge [loop left] node [left] {$1_A$} (m-1-1)
		(m-1-2) edge [loop right] node [right] {$1_B$} (m-1-2)
		;
		\end{tikzpicture}
		\]
		\item Ad\'{a}mek, Herrlich, and Strecker (2004) \cite{agat:AHS2004} posited further that $\text{Hom}(A,B) \in \text{Mor}\mathbf{A}$ pairwise disjoint (i.e. $\text{Hom}(A,B) \bigcap \text{Hom}(C,D) \neq \emptyset$ if $C\neq A$ or $D\neq B$)
	\end{enumerate}
\end{enumerate}	
\end{definition}

\subsection{Examples of categories}

\begin{itemize}
	\item $\textbf{Set} = (\text{Obj}{(\textbf{Set})}, \text{Hom}{\textbf{Set}},\mathbf{1},\circ)$ where \\
	$\text{Obj}{(\textbf{Set})}$ is the class of all sets \\
	$\text{Hom}{\textbf{Set}}$ is the class of all functions on a set to another set
	\item $\textbf{Vec}$
	
	\[
	\begin{aligned}
	& \text{Obj}\textbf{Vec} & \equiv \text{ all real vector spaces } \\ 
	& \text{Mor}\textbf{Vec} & \equiv \text{ all linear transformations between them (between real vector spaces) }
	\end{aligned}
	\]
	\item \emph{preorder} cf. Ch. 1, Example 7, pp. 8, Awodey (2010) \cite{agat:Awod2010}. \\
	\emph{preorder} $P$, set $P$ equipped with binary relation $p\leq q $, that's both \\
	reflexive: $a\leq a$, and \\
	transitive: if $a \leq b$ and $b \leq c$, then $a\leq c$. 
	
	\quad \\
	$\forall \, $ preorder $P$ can be regarded as a category: \\
	$\text{Obj}(P) = P$ \\
	$a \to b$ iff $a\leq b \in \text{Mor}P$. \\
	identity: $1 :a \to a$ since $a\leq a$ \\
	composition: $f:a\to b$, $g:b\to c$, $(g\circ f)a = g(b) = c$, $a\leq b$ and $b\leq c$, so $a\leq c$. 
	
	Going in the other direction, a category with at most 1 arrow between 2 objects determines a preorder.
	\item \textbf{Monoid}.  Consider a monoid as a triple $(M, \cdot, e)$.  \\
	Every semigroup $(M,\cdot)$ (recall that a \emph{semigroup} is a set $S$ with binary operation $\cdot $, i.e. s.t.
	
	\quad \,  $\begin{aligned} & \quad \\
	& S\times S \xrightarrow{\cdot } S \\ 
	& \forall \, a,b ,c \in S, \, (a\cdot b)\cdot c = a\cdot (b\cdot c) \quad \, \text{ (associativity) } \end{aligned}$ 
	
	\quad \, (but no inverse, necessarily!)) that also has a unit $e$ can be made into a category $\mathbf{C}$ 
	
	$\Longrightarrow \mathbf{C}(M,\cdot ,e) = (\text{Obj}(\mathbf{C}), \text{Hom}(\mathbf{C}), \mathbf{1}, \circ)$, a category $\mathbf{C}$ with only 1 object, i.e. $\text{Obj}(\mathbf{C}) = \lbrace M \rbrace$, so that \\
	$\text{Obj}(\mathbf{C}) = \lbrace M \rbrace$ \\
	$\text{Hom}(M,M) = M$ \\
	$\mathbf{1}_M = e$ \\
	$y \circ x = y \cdot x$
	
	cf. pp. 10, Example 12 of Awodey (2010) \cite{agat:Awod2010}.

\subsubsection{Monoid as a category}

Recall the definition of a \emph{category} $\mathbf{A}$: category $\mathbf{A} = (\text{Obj}(\mathbf{A}), \text{Mor}{\mathbf{A}}, 1, \circ)$ s.t.\\
$A \in \text{Obj}\mathbf{A}$, \\
$\forall \, A, B \in \text{Obj}\mathbf{A}$, $\exists \, \text{Hom}(A,B) \subseteq \text{Mor}\mathbf{A}$, \\
$\forall \, f \in \text{Hom}(A,B)$, $f:A \to B$ is a morphism.\\
$\forall \, A \in \text{Obj}\mathbf{A}$, $\exists \, 1_A : A \to A$, i.e. $\exists \, 1_A \in \text{Hom}_{\mathbf{A}}(A,A)$. \\
$\forall \, A, B, C \in \text{Obj}\mathbf{A}$,  \\
\[
\begin{aligned} 
	& \text{Hom}_{\mathbf{A}}(A, B) \times \text{Hom}_{\mathbf{A}}(B,C) \to \text{Hom}_{\mathbf{A}}(A, C) \\ 
	& (f,g) \mapsto g\circ f \end{aligned} \] \\

Recall the definition of a \emph{monoid}:
monoid $(M, \cdot, e)$, set $M$, \\
$\cdot : M\times M \to M$, \\
$e \in M$ s.t. $\forall \, x \in M$, $e\cdot x = x \cdot e = x$ \\

$\Longrightarrow \mathbf{C}(M)$ (monoid as a category) \\
$\text{Obj}\mathbf{C}(M) = \lbrace M \rbrace$ \\
$\text{Hom}(M, M) = \text{Mor}\mathbf{C}(M)$ and $\forall \, f \in \text{Hom}(M, M)$, 
\[
f = m \in M
\]
i.e. the morphism \emph{is} an element in monoid $M$ (there's a 1-to-1 correspondence).

$\forall \, n \in M$, $f:n \mapsto m\cdot n \in M$. \\
$1_M : M \to M$ \\
$1_M : x \mapsto e\cdot x = x$ \\
$m\cdot n \cdot x = m\cdot (n\cdot x) = (m\cdot n) \cdot x$ (composition)


	
\subsubsection{Examples of monoids}
	Examples of monoids: $\mathbb{N}, \mathbb{Q}, \mathbb{R}$ under either addition and $e=0$, or multiplication and $e=1$. \\
	$\forall \, $ set $X$, set of functions from $X$ to $X$, $\text{Hom}_{\mathbf{Sets}}(X, X)$ is a monoid under composition. 
	
	In general, \\
	$\forall \, $ object $C \in \text{Obj}(\mathbf{C})$, $\forall \, $ category $\mathbf{C}$, $\text{Hom}_{\mathbf{C}}(C,C)$ is a monoid under composition of $\mathbf{C}$. 
	
	Since monoids are structured sets, $\exists \, $ category $\textbf{Mon}$, s.t. $\text{Obj}\mathbf{Mon} \ni $ monoids, $\text{Mor} \textbf{Mon} \ni $ functions that preserve monoid structure.
	
	In detail, homomorphism from monoid $M$ to monoid $N$ is function $h:M \to N$ s.t. $\forall \, m,n \in M$. 
	
	\[
	h(m \cdot_M n) = h(m) \cdot_N h(n)
	\]
	and
	\[
	h(e_M) =  e_N
	\]
	
	\quad \\ 
	Check: Consider functor $F: M \to N$. \\
	$F: \text{Hom}_M(a,b) \to \text{Hom}_N(F(a), F(b))$
	\[
	\begin{gathered}
	m\cdot a =b \qquad \, n \cdot_N F(a)  = F(b) \\
	F(m) \cdot_N F(a) = F(b) = F(m \cdot_M a) 
	\end{gathered}
	\]
	so functor $F$ is a monoid homomorphism, and so \\
	\emph{a monoid homomorphism from $M$ to $N$ is the same thing as a functor from $M$ regarded as a category to $N$ regarded as a category.}
	\item \textbf{Poset} cf. Ch. 1, pp. 6, Awodey (2010) \cite{agat:Awod2010}. partially ordered set or poset is a set $A$, equipped with binary relation $a \leq_A b$ s.t. $\forall \, a, b, c \in A$, \\
	reflexivity: $a \leq_A a$ \\
	transitivity: if $a\leq_A b$, $b\leq_A c$, then $a\leq_A  c$ \\
	antisymmetry: if $a\leq_A b$, $b\leq_A a$, then $a=b$.
	
	An arrow from a poset $A$ to a poset $B$ is a function,
	\[
	m : A \to B
	\]
that is monotone, i.e. $\forall \, a, a' \in A$, 
\[
\text{ if } a \leq_A a', \text{ then } m(a) \leq_B m(a')
\]	
category $\textbf{Pos}$ of posets, and monotone functions, 
\[
\forall \, A \in \text{Obj}(\textbf{Pos}), \, \exists \, 1_A : A \to A, \text{ since } 
\]
since if $a \leq_A a'$, then $a\leq_A a'$, so $1_A$ monotone.

\quad \\ 
\emph{composition}: $\forall \, f:A \to B, \, g:B \to C$, \, $\forall \, A,B,C \in \text{Obj}(\textbf{Pos})$, $f,g$ monotone, \\
$\forall \, a \leq_A a'$, $f(a) \leq_B f(a')$, and since \\
$\forall \, b \leq_B b'$, $g(b) \leq_C g(b')$, then $g(f(a)) \leq_C g(f(a'))$, or $(g\circ f)(a) \leq_C (g\circ f)(a') \, \forall \, a\leq_A a'$, so $g\circ f : A \to C$ is monotone.

$\text{Mor}\textbf{Pos} = \bigcup_{A, B \in \text{Obj}(\textbf{Pos}) } \text{ monotone functions $A \to B$ }$  

\quad \\ 
A \emph{poset} is a \emph{preorder}, satisfying additional condition of antisymmetry:
if $a\leq b$, and $b\leq a$, $a=b$. \\
e.g. $\forall \, X$, power set $P(X)$ is a poset under inclusion $U\subseteq V$, $\forall \, $ subsets $U,V \subseteq X$.

\begin{itemize}
	\item cf. Ch. 1, Example 8, pp. 9, Awodey (2010) \cite{agat:Awod2010}.
	$\textbf{Pos}$ categories and functors: functor $F: \mathbf{D} \to \mathbf{Q}$ between poset categories $\mathbf{P}, \mathbf{Q}$. \\
	They are the monotone functions.
	\begin{proof} 
	Consider 
	\[
	F : \text{Hom}_{\mathbf{P}}(A,B) \to \text{Hom}_{\mathbf{Q}}(F(A), F(B))
	\]
	$f\in \text{Hom}_{\mathbf{P}}(A,B)$ so if $a\leq_A a'$, $f(a) \leq_B f(a')$
	
	Let $g\in \text{Hom}_{\mathbf{Q}}(F(A), F(B))$, so if $c\leq_{F(A)} c'$, $g(c)\leq_{F(B)} g(c')$. \\
	
	$F(f) \in \text{Hom}_{\mathbf{Q}}(F(A), F(B))$ (by definition of a functor), so if $c\leq_{F(A)} c'$, $F(f)(c) \leq_{F(B)} F(f)(c')$ or $F(f(c)) \leq_{F(B)} F(f(c'))$. \\
	
	So $F$ itself is monotone.
	\end{proof}
\end{itemize}
\item cf. \emph{An example from computer science}, Example 10, pp. 9 of Awodey (2010) \cite{agat:Awod2010}.
Given a functional programming language $L$, $\exists \, $ associated category, \\
$\text{Obj}\mathbf{L} \ni $ data types of $L$ \\
$\text{Mor}\mathbf{L} \ni $ computable functions of $L$ ("processes", "procedures", "programs") \\
composition (of 2 such programs) $X \xrightarrow{f} Y \xrightarrow{g} Z$ given by applying $g$ to output of $f$ \\
(notation note: $g\circ f = f; g$) \\
identity is the "do nothing" program. 

\quad \\
if $\mathbf{C}(L)$ is the category just defined, then \\
denotational semantics of language $L$ in category $\mathbf{D}$ of say Scott demands is simply functor
\[
S:\mathbf{C}(L) \to \mathbf{D} 
\]
since $S$ assigns domains to the types of $L$, \\
continuous functions to programs.

This example and Example 9, pp. 9, of Awodey (2010) \cite{agat:Awod2010} are related to the notion of "Cartesian closed category."

From Barr and Wells (2012) \cite{agat:BaWe2012}, a \emph{functional programming language} $L$ has
\begin{enumerate}
	\item primitive data types (built into the language) 
	\item constants of each type
	\item operations, which are functions between types
	\item constructors, which are applied to data types and operations, to produce derived data types and operations of the language.	
\end{enumerate}

	For a $\mathbf{C}(L)$ category corresponding to functional programming language $L$, \\
	
	$\mathbf{C}(L)$ category corresponding to functional programming language $L$, \\
	$\text{Obj}\mathbf{C}(L) = \lbrace \text{ types of $L$ } \rbrace$ \\
	$\text{Mor}\mathbf{C}(L) = \lbrace \text{ operations (functions; primitive and derived) of $L$ } \rbrace$ \\

identity: $\forall \, $ type $T \in \text{Obj}\mathbf{C}(L)$, $\exists \, 1_T :T \to T$ i.e. $1_T \in \text{Hom}(T,T)$. \\

composition: $\forall \, $ types $T,U,V \in \text{Obj}\mathbf{C}(L)$, defines composition 
\[
\begin{aligned}
\text{Hom}(T,U) \times \text{Hom}(U,V) & \to \text{Hom}(T,V) \\ 
(f,g) & \mapsto g \circ f
\end{aligned}
\]

Assume $L$ has a do nothing operation $1_T$, and composition constructor, $\text{Hom}(T,U) \times \text{Hom}(U,V) \to \text{Hom}(T, V)$. \\

Add additional type $1$ s.t. $\forall \, $ type $T \in \text{Obj}\mathbf{C}(L)$, $\exists \, !$ operation (function) to $1, T \to 1$ \\
Interpret each constant $c$ of type $A$ as arrow (morphism) $c:1\to A$.  \\

As a concrete example (cf. Example 2.2.5 of Barr and Wells (2012) \cite{agat:BaWe2012}), suppose a simple language with 3 data types, \texttt{int}, \texttt{bool}, \texttt{char}. \\

For \texttt{int}, $\exists \, $ constant $0:1 \to \texttt{int}$, consider $\texttt{succ}:\texttt{int} \to \texttt{int} \in \text{Hom}(\texttt{int}, \texttt{int})$, \\
For \texttt{bool}, $\exists \, $ constants $\texttt{true} : 1 \to \texttt{bool}$, $\texttt{false} : 1 \to \texttt{bool}$; consider $\neg \in \text{Hom}(\texttt{bool}, \texttt{bool})$, s.t.
\[
\begin{aligned}
& \neg \circ \texttt{true} = \texttt{false} \\ 
& \neg \circ \texttt{false} = \texttt{true} 
\end{aligned}
\]

For \texttt{char}, $\exists \, $ constants $c:1 \to \texttt{char}$, \, $\forall \, \texttt{char} c$.

Since $\forall \, T,U \in \text{Obj}(\mathbf{C}(L))$, $\exists \, \text{Hom}(T,U) \subseteq \text{Mor}\mathbf{C}(L)$. So consider $\texttt{ord}:\texttt{char} \to \texttt{int}$, $\texttt{chr}: \texttt{int} \to \texttt{char}$, s.t. 
\[
\texttt{char} \circ \texttt{ord} = 1_{\texttt{char}}
\]

So $\text{Obj}\mathbf{C}(L) = \lbrace \texttt{int}, \texttt{bool}, \texttt{char}, 1 \rbrace$, \\
\phantom{So } $\text{Mor}\mathbf{C}(L)$ consists of all programs.

See 5.3.14 Record Types and 5.7.6, 14.2 for flow of control in Barr and Wells (2012) \cite{agat:BaWe2012} to complete the programming language. For here, see section \ref{agat:SubSubSec:RecordTypesProduct}

\item \textbf{Groups}, cf. pp. 11, Def. 1.4, Sec. 1.5 \textbf{Isomorphisms}, Awodey (2010) \cite{agat:Awod2010}.

A \emph{group} $G$ is a monoid with inverse $g^{-1}$, $\forall \, g \in G$. \\
Thus, $G$ is a category with 1 object, $\forall \, f \in \text{Mor}G$ is an isomorphism.

$\forall \, $ set $X$, group $\text{Aut}(X)$ of automoprhisms (or "permutations") of $X$ i.e. $f:X\to X$ \\
$\text{Aut}(X)$ is closed under composition $\circ$ because a permutation of a permutation is another permutation.

Homomorphisms of groups $h: G\to H$ is just a homomorphism of monoids, preserving inverses.\\
Given $h$, then $h(gg') = h(g) h(g')$. \\
$G,H$ are monoids. Then $h$ monoid homomorphism.
\end{itemize}

Concrete categories: informally, categories which objects are sets, possibly equipped with some structure, and arrows are certain, possibly structure-preserving functions.

From pp. 10 of Mac Lane (1978) \cite{agat:MacL1978}, Leinster (2014) \cite{agat:Lein2014}, Examples 1.1.8 (Categories as mathematical structures), pp. 13, and Sec. 1.4 Examples of Categories, 5. Finite Categories, pp.7 of Awodey (2010) \cite{agat:Awod2010}.

Category $\mathbf{1}$.  $\text{Obj}{\mathbf{1}} = \lbrace A \rbrace$ \\
$\text{Mor}{\mathbf{1}} = \lbrace 1_A \rbrace$; $1_A : A \to A$.  \\

The only composition to consider is $1_A \circ 1_A : A \to A$, and $1_A \circ 1_A \circ 1_A$, so on and so for.

Category $\mathbf{2}$.  $\text{Obj}{\mathbf{2}} = \lbrace A, B \rbrace$. \\
$\text{Hom}(A, B) = \lbrace f \rbrace$, $f:A\to B$. \\

The only "non-trivial" composition to consider is $1_B \circ f \circ 1_A : A \to B$. \\
 
Category $\mathbf{3}$. pp. 11, Sec. 2 Categories of Mac Lane (1978) \cite{agat:MacL1978}, Example 5. "Finite Categories", Sec. 1.4, pp. 7 of Awodey (2010) \cite{agat:Awod2010}. 

$\text{Obj}{\mathbf{3}} = \lbrace A, B, C \rbrace$ \\
$\text{Hom}{(A, B)} = \lbrace f \rbrace$ \\
$\text{Hom}{(A, C)} = \lbrace g \rbrace$ \\
$\text{Hom}{(B, C)} = \lbrace h \rbrace$ \\
$\text{Hom}{(C, \cdot)} = \emptyset $ \\
$\text{Hom}{(B, A)} = \emptyset$ \\

The only "non-trivial" compositions to consider are the following:
$1_C \circ h \circ f : A \to C$, \\
$1_C \circ g$, $1_C \circ h$ (and likewise) \\


\subsection{Duality, opposite category}

Given a category $\mathbf{A} = ( \text{Ob}, \text{hom}_{\mathbf{A}}, 1, \circ)$, 
\begin{definition}[dual opposite category]
	\textbf{dual} or \textbf{opposite} category of $\mathbf{A} = (\text{Obj}(\mathbf{A}), \text{Mor}\mathbf{A}, \mathbf{1}, \circ)$, denoted $\mathbf{A}^{\text{op}}$, is   
	\begin{equation}
	\mathbf{A}^{\text{op}}  =  (\text{Obj}(\mathbf{A}), \text{Mor}\mathbf{A}^{\text{op}}, \mathbf{1}, \circ^{\text{op}})
	\end{equation}
	s.t.
\begin{itemize}
	\item 
	\begin{equation}
	\text{Obj}(\mathbf{A}^{\text{op}}) = \text{Obj}(\mathbf{A})
	\end{equation}
	\item $\forall \, A, B \in \text{Obj}(\mathbf{A}^{\text{op}})$, $\text{Hom}_{\mathbf{A}^{\text{op}}}(A,B) \subseteq \text{Mor}\mathbf{A}^{\text{op}}$, 
		\begin{equation}
\text{Hom}_{\mathbf{A}^{\text{op}}}(A,B) = \text{Hom}_{\mathbf{A}}(B,A) \subseteq \text{Mor}\mathbf{A}
		\end{equation}
		\item \emph{Define} the new composition 
		\begin{equation}
\begin{gathered}
		f \circ^{\text{op}} g \text{ of } \begin{aligned} & \quad \\
		& g \in \text{Hom}_{\mathbf{A}^{\text{op}}}(C,B) \\
		& f \in \text{Hom}_{\mathbf{A}^{\text{op}}}(B,A) 
		\end{aligned}		 \\
		\text{ then } \\
		f \circ^{\text{op}} g = g\circ f
\end{gathered}
		\end{equation}
\[
\begin{tikzpicture}
\matrix (m) [matrix of math nodes, row sep=3em, column sep=4em, minimum width=2em]
{
	A & B & C \\
};
\path[->]
(m-1-1) edge node [above] {$f$} (m-1-2)
(m-1-2) edge node [above] {$g$} (m-1-3);
		\path[-{>[scale=1.15]}, line width=0.5pt]
(m-1-1) edge[bend right=45] node [below] {$g\circ f$} (m-1-3);
\end{tikzpicture} \qquad \qquad \, \, \text{ dual :} \begin{tikzpicture}
\matrix (m) [matrix of math nodes, row sep=3em, column sep=4em, minimum width=2em]
{
	A & B & C \\
};
\path[->]
(m-1-3) edge node [above] {$g$} (m-1-2)
(m-1-2) edge node [above] {$f$} (m-1-1);
		\path[-{>[scale=1.15]}, line width=0.5pt]
(m-1-3) edge[bend left=45] node [below] {$f\circ^{\text{op}} g$} (m-1-1);
\end{tikzpicture}
\]		
or, equivalently (notation-wise)
\[
\begin{tikzpicture}
\matrix (m) [matrix of math nodes, row sep=3em, column sep=4em, minimum width=2em]
{
	A & B & C \\
};
\path[->]
(m-1-3) edge node [above] {$g$} (m-1-2)
(m-1-2) edge node [above] {$f$} (m-1-1);
		\path[-{>[scale=1.15]}, line width=0.5pt]
(m-1-3) edge[bend left=45] node [below] {$f \circ g$} (m-1-1);
\end{tikzpicture}
\qquad \qquad \, \, \text{ dual :} 
\begin{tikzpicture}
\matrix (m) [matrix of math nodes, row sep=3em, column sep=4em, minimum width=2em]
{
	A & B & C \\
};
\path[->]
(m-1-1) edge node [above] {$f$} (m-1-2)
(m-1-2) edge node [above] {$g$} (m-1-3);
		\path[-{>[scale=1.15]}, line width=0.5pt]
(m-1-1) edge[bend right=45] node [below] {$g\circ^{\text{op}} f$} (m-1-3);
\end{tikzpicture} 
\]		
in that 
\[
\begin{gathered}
g \circ^{\text{op}} f \text{ of } \begin{aligned} & \quad \\
& f \in \text{Hom}_{\mathbf{A}^{\text{op}}}(A,B) \\
& g \in \text{Hom}_{\mathbf{A}^{\text{op}}}(B,C) 
\end{aligned}		 \\
\text{ then } \\
g \circ^{\text{op}} f = f\circ g
\end{gathered}
\]
\end{itemize}
\end{definition}

i.e. (in summary) $\mathbf{C}^{\text{op}}$ s.t. $\text{Obj}\mathbf{C}^{\text{op}} = \text{Obj}\mathbf{C}$, \\
Denote $\forall \, \overline{C} \in \text{Obj}\mathbf{C}^{\text{op}}$, $\overline{C} = C \in \text{Obj}\mathbf{C}$. \\
$\text{Hom}_{\mathbf{C}^{\text{op}}}(\overline{A}, \overline{B}) \ni \overline{f} : \overline{C} \to \overline{D}$, in $\mathbf{C}^{\text{op}}$, for $f:D\to C$ in $\mathbf{C}$.  \\

\[
\begin{gathered}
1_{\overline{C}} = 1_C \\
	\overline{f} \circ \overline{g} \equiv \overline{f} \circ_{\text{op}} \overline{g} = g\circ f 
\end{gathered}
\]

Diagram in $\mathbf{C}$:

\[
\begin{tikzpicture}
\matrix(m)[matrix of math nodes, row sep=4em, column sep=4em]
{
	A   &  B \\
	  & C  \\
};
\path[->]
(m-1-1) edge node[above]{$f$} (m-1-2)
edge node[left]{$g\circ f$} (m-2-2) 
(m-1-2) edge node[right]{$g$} (m-2-2);
\end{tikzpicture} 
\]

Diagram in $\mathbf{C}^{\text{op}}$:

\[
\begin{tikzpicture}
\matrix(m)[matrix of math nodes, row sep=4em, column sep=4em]
{
	\overline{A}   &  \overline{B} \\
	& \overline{C}  \\
};
\path[->]
(m-1-2) edge node[above]{$\overline{f}$} (m-1-1)
(m-2-2) edge node[left]{$\overline{f}\circ \overline{g}$} (m-1-1)
edge node[right]{$\overline{g}$} (m-1-2);
\end{tikzpicture} 
\]

e.g. $\mathbf{Sets}$ is dual to category of complete, atomic Boolean algebras.

e.g. if $\mathbf{A} = (M,\cdot, e)$ monoid, then $\mathbf{A}^{\text{op}} = (M, \widehat{\cdot},e)$ where $a\widehat{\cdot} b = b\cdot a$


\subsubsection{Example}

\begin{itemize}
	\item $\text{Vec}^{\text{op}}$
	\[
	\textbf{Vec}^{\text{op}} = (\text{Obj}(\textbf{Vec}), \text{Hom}_{\textbf{Vec}^{\text{op}}}, 1, \circ^{\text{op}})
	\]
	s.t.
	\[
	\text{Hom}_{\textbf{Vec}^{\text{op}} }(W,V) = \text{Hom}_{\textbf{Vec}}(V,W)
	\]
\[	
	\begin{tikzpicture}
	\matrix (m) [matrix of math nodes, row sep=3em, column sep=4em, minimum width=2em]
	{
		U & V & W \\
	};
	\path[->]
	(m-1-1) edge node [above] {$f$} (m-1-2)
	(m-1-2) edge node [above] {$g$} (m-1-3);
		\path[-{>[scale=1.15]}, line width=0.5pt]
	(m-1-1) edge[bend right=45] node [below] {$g\circ f$} (m-1-3);
	\end{tikzpicture} \qquad \, \begin{tikzpicture}
	\matrix (m) [matrix of math nodes, row sep=3em, column sep=4em, minimum width=2em]
	{
		U & V & W \\
	};
	\path[->]
	(m-1-3) edge node [above] {$g$} (m-1-2)
	(m-1-2) edge node [above] {$f$} (m-1-1);
		\path[-{>[scale=1.15]}, line width=0.5pt]
(m-1-3) 	edge[bend left=45] node [below] {$f\circ^{\text{op}} g$} (m-1-1);
	\end{tikzpicture}
\]
\end{itemize}

\subsection{Kinds of morphisms}

cf. pp. 11, Definition 1.3 of Awodey (2010) \cite{agat:Awod2010}
\begin{definition}[isomorphism]
	\textbf{isomorphism} - $\forall \, $ category $\mathbf{C}$, $\text{Mor}\mathbf{C} \ni $ morphism $f:A \to B$ is an isomorphism if $\exists \, g : B \to A$ s.t. $f\circ g = 1_B$, $g\circ f = 1_A$, $g$ unique. $g$ called inverse of $f,f^{-1}$
	\[
	\begin{gathered}
	\begin{tikzpicture}
\matrix (m) [matrix of math nodes, row sep=3em, column sep=4em, minimum width=2em]
{
	A & B \\
};
\path[->]
(m-1-1) edge node [above] {$f$} (m-1-2)
		edge [loop left] node [left] {$1_A$} (m-1-1);
		\path[-{>[scale=1.15]}, line width=0.5pt]
(m-1-2) edge[bend left=45] node [below] {$g$} (m-1-1);
\end{tikzpicture}	\qquad \, 
	\begin{tikzpicture}
\matrix (m) [matrix of math nodes, row sep=3em, column sep=4em, minimum width=2em]
{
	B & A \\
};
\path[->]
(m-1-1) edge node [above] {$g$} (m-1-2)
		edge [loop left] node [left] {$1_B$} (m-1-1);
\path[-{>[scale=1.15]}, line width=0.5pt]
(m-1-2) edge[bend left=45] node [below] {$f$} (m-1-1);
\end{tikzpicture}	
	\end{gathered}	
	\]
\end{definition}

\begin{theorem}[Cayley]
cf. pp. 11, Sec. 1.5. \textbf{Isomorphisms}, Awodey (2010) \cite{agat:Awod2010}

$\forall \, $ group $G$, 
\[
G \cong \text{ group of permutations }
\]
(isomorphic $\equiv \cong$)

\end{theorem} 

\begin{proof}
	Define Cayley representation $\overline{G}$ of $G$ to be following group of permutations: \\
	underlying set of $\overline{G}$ is just $G$, \\
	$\forall \, g\in G$, we have permutation $\overline{g}$, defined $\forall \, h \in G$, 
	\[
	\overline{g}(h) = g\cdot h
	\]
	Let $\overline{g} = \overline{h}$. \\
	$\forall \, g' \in G$, \\
	$\begin{aligned}
	& \overline{g}(g') = gg' \\ 
	& \overline{h}(g') = hg'
	\end{aligned}$ \\
	$g' \in G$, so $\exists \, (g')^{-1}$ s.t. $g'(g')^{-1} =e$, so $gg'(g')^{-1}= g=h$.
	
	\quad \\
	Define homomorphisms $\begin{aligned} & \quad \\ 
	& i : G \to \overline{G} \\
	& i(g) = \overline{g}
	\end{aligned}$, \qquad \, $\begin{aligned} & \quad \\ 
	& j : \overline{G} \to G \\
	& j(\overline{g}) =  g
	\end{aligned}$
	
	\[
	\begin{gathered} 
	i(gg') = \overline{(gg')} \\
	\begin{aligned}
	& \overline{(gg')} = h = gg' h = g(g'h) = g(\overline{g}'(h)) = \overline{g}(\overline{g}'(h)) = i(g) = i(g') h \\ 
	& \Longrightarrow i(gg') = i(g) \circ i(g') \\
	& j(\overline{g}) = j(\overline{g}') h = j(\overline{g}) (g'h) = gg'h = j(\overline{gg'})h = j(\overline{g} \cdot \overline{g}') h \\
	& j(\overline{g}) \cdot j(\overline{g}') = j(\overline{g} \cdot \overline{g}')
	\end{aligned} \\
	\begin{aligned}
	& i\circ j = 1_{\overline{G}}, \, j \circ i = 1_G \text{ since } \\
	& i \circ j (\overline{g}) = i\circ g = \overline{g} \\
	& j\circ i (g) = j(\overline{g}) = g
	\end{aligned}
\end{gathered} 	
	\]
	$i,j$ are isomorphisms between $G,\overline{G}$, which is in category $\textbf{Groups}$ of groups and group homomorphisms. \\
	permutations $\overline{g} \in \overline{G}$ are themselves isomorphisms in $\textbf{Sets}$. \\
	Cayley's Thm. says that any abstract group can be represented as a "concrete" one, i.e. group of permutations of a set.
\end{proof}

Example 1.1.5 (Leinster (2014) \cite{agat:Lein2014}) isomorphisms in $\textbf{Set}$ are exactly bijections. \\
$\Longrightarrow $ function has 2-sided inverse iff function is injective and surjective (this is not trivial)

Example 1.1.7 Leinster (2014) \cite{agat:Lein2014} isomorphisms in $\textbf{Top}$ are exactly homeomorphisms. \\
bijective map in $\textbf{Top}$ is not necessarily an isomorphism:
\[
\begin{aligned}
	[0, 1) & \to \lbrace z \in \mathbb{C} | |z| = 1 \rbrace \\
	t & \mapsto e^{2\pi i t}
\end{aligned}
\]
is a continuous bijection but not a homeomorphism (cont. bijection with cont. $f^{-1}$). 
\[
\ln{ \left( \frac{w}{2\pi i} \right) } = f^{-1}(w)
\]

Example 1.1.8 Leinster (2014) \cite{agat:Lein2014} (Categories as mathematical structures) (b) \textbf{Discrete} categories contain no maps at all apart from identities. Just a class of objects.

\exercisehead{1.1.13 pp. 16, Leinster (2014) \cite{agat:Lein2014}}

Suppose for morphism $f: A \to B$, $\exists \,$ inverse $g$, s.t. $f\circ g = 1_B$, $g\circ f = 1_A$ \\
Now $f^{-1} \circ f = 1_A$, $f\circ f^{-1} = 1_B$ \\
\phantom{Now} $f^{-1} \circ f \circ g = f^{-1} \circ 1_B = 1_A \circ g = f^{-1} \circ 1_B \Longrightarrow g= f^{_1}$ 


\begin{definition}[endomorphism]
	\textbf{endomorphism} - morphism with same source and target, that is, morphism $f:A \to A$
\end{definition}

\begin{definition}[automorphism]
\textbf{automorphism} - endomorphism which is an isomorphism
\end{definition}

\begin{definition}[parallel]
	\textbf{parallel} - 2 morphisms $f,g$ are parallel if they have same source and same target:
	\[
	\begin{aligned}
	& f:A \to B \\
	& g:A \to B
	\end{aligned}
	\]
\end{definition}

\begin{definition}[monomorphism]
	\text{monomorphism} - morphism $f: A \to B$ is a monomorphism if $\forall \, $ pair of parallel $\begin{aligned}
	& \quad \\
	& g_1: C \to A \\
	& g_2:C \to A
	\end{aligned}$, 
	\begin{equation}
	f\circ g_1 = f\circ g_2 \text{ implies } g_1 = g_2
	\end{equation}
	i.e.
	\[
	\begin{tikzpicture}
	\matrix (m) [matrix of math nodes, row sep=3em, column sep=4em, minimum width=2em]
	{
		C & B \\
	};
	\path[->]
	(m-1-1) edge node [above] {$f \circ g_1$} (m-1-2);
	\path[-{>[scale=1.15]}, line width=0.5pt]
	(m-1-2) edge[bend left=45] node [below] {$f\circ g_2$} (m-1-1);
	\end{tikzpicture} \text{ implies } C \xrightarrow{ g_1 = g_2 } A	
	\]
\end{definition}

\begin{definition}[epimorphism]
\textbf{epimorphism} - morphism $f:A \to B$ is an epimorphism if $f^{\text{op}} : B^{\text{op}} \to A^{\text{op}}$ is a monomorphism in $\mathbf{A}^{\text{op}}$. \\
Hence $f$ epimorphism iff $\forall \, $ parallel morphisms $\begin{aligned} & \quad \\ 
	& g_1 : B \to C \\
	& g_2 : B\to C\end{aligned}$, $g_1 \circ f = g_2 \circ f $ \\
	 implies $g_1 = g_2$
\end{definition}

\begin{proposition}[monomorphism, epimorphism iff injective]
	$f$ monomorphism iff $f \circ : \text{Hom}_{\mathbf{A}}(C,A) \to \text{Hom}_{\mathbf{A}}(C,B)$ injective \, $\forall \, C \in \text{Obj}(\mathbf{A})$, i.e.
	\begin{equation}
	\begin{gathered}
\begin{gathered} 
	\text{Hom}_{\mathbf{A}}(C, A) \xrightarrow{f \circ } \text{Hom}_{\mathbf{A}}(C,B) \\
 g_1 , g_2 \xmapsto{f \circ} f\circ g_1 , f\circ g_2
\end{gathered}, \\ \text{ then } \\
f \circ \text{ is injective if } \\
	f \circ g_1 = f \circ g_2 \Longrightarrow g_1 = g_2
\end{gathered} 
	\end{equation}
$f$ epimorphism iff map $\circ f : \text{Hom}_{\mathbf{A}}(B,C) \to \text{Hom}_{\mathbf{A}}(A,C)$ injective $\forall \, C \in \text{Obj}(\mathbf{A})$
\begin{equation}
\begin{gathered}
\begin{gathered} 
\text{Hom}_{\mathbf{A}}(B, C) \xrightarrow{\circ f} \text{Hom}_{\mathbf{A}}(A,C) \\
g_1 , g_2  \xmapsto{ \circ f} g_1\circ f , g_2 \circ f
\end{gathered}, \\ \text{ then } \\
\circ f \text{ is injective if } \\ 
g_1 \circ f = g_2 \circ f \Longrightarrow g_1 = g_2
\end{gathered}
\end{equation}
\end{proposition}

\begin{definition}[inverses]
	$\forall \, $ 2 morphisms, $f:X\to Y$, $g: Y \to X$ s.t. $f\circ g= 1_Y$,\\
	$f$ is called left inverse of $g$, $g$ is called right inverse of $f$.
	
	We also say, $g$ is a section of $f$, or $f$ is a cosection of $g$. \\
	$f$ is an epimorphism, $g$ is a monomorphism.
\end{definition}

\subsection{More definitions with categories}

\begin{definition}[subcategory]
category $\mathbf{A}'$, $\mathbf{A}' \subset \mathbf{A}$, if $\text{Obj}(\mathbf{A}') \subset \text{Obj}(\mathbf{A})$, $\text{Hom}_{\mathbf{A}'}(A,B) \subset \text{Hom}_{\mathbf{A}}(A,B)$, $\forall \, A, B \in \mathbf{A}'$.

Composition in $\mathbf{A}'$ is induced by composition in $\mathbf{A}$.\\
identity morphisms in $\mathbf{A}'$ are identity morphisms in $\mathbf{A}$
\end{definition}

\begin{definition}[full subcategory]
	subcategory $\mathbf{A}'$ of $\mathbf{A}$ is full if $\text{Hom}_{\mathbf{A}'}(A,B) = \text{Hom}_{\mathbf{A}}(A,B)$, $\forall \, A, B \in \mathbf{A}'$	
\end{definition}

\begin{definition}[saturated subcategory]
full subcategory $\mathbf{A}'$ of $\mathbf{A}$ saturated if $A \in \mathbf{A}$ belongs to $\mathbf{A}'$ whenever $A$ is isomorphic to object of $\mathbf{A}'$	
\end{definition}

\begin{definition}[discrete category]
	\textbf{discrete} - discrete category if all morphisms are identity morphisms.
\end{definition}

\begin{definition}[nonempty category]
	\textbf{nonempty} - nonempty category if $\text{Obj}(\mathbf{A})$ is nonempty
\end{definition}

\begin{definition}[groupoid]
\textbf{groupoid} - category $\mathbf{A}$ is a \textbf{groupoid} if all morphisms are isomorphisms.
\end{definition}

\begin{definition}[finite category]
	finite - finite category if set of all morphisms in $\mathbf{A}$ (hence, in particular, set of objects) is a finite set
\end{definition}

\begin{definition}[connected]
	connected category $\mathbf{A}$ if it's nonempty, and $\forall \, A,B \in \text{Obj}\mathbf{A}$, $\exists \, $ finite sequence of objects $(A_0 \dots A_n)$, $A_0 = A$, $A_n = B$, s.t. at least 1 of the sets $\text{Hom}_{\mathbf{A}}(A_j, A_{j+1})$ or $\text{Hom}_{\mathbf{A}}(A_{j+1}, A_j)$ is nonempty $\forall \, j \in \mathbb{N}$, with $0\leq j \leq n-1$
\end{definition}

\begin{definition}[monoid $M$]
	\textbf{monoid} $M$ (set endowed with internal product with associative and unital law) is nothing but a category with only 1 object (to $M$, associate category $\mathbf{M}$, with single object $A$, and morphisms $\text{Hom}_{\mathbf{M}}(A,A) = M$)
\end{definition}

cf. Def. 1.2.5 of Kashiwara and Schapira (2006) \cite{agat:KaSch2006}.

\begin{definition}[Morphisms as a category]
	Let category $\mathcal{C} \equiv \mathbf{A}$.
	
	$\text{Mor}(\mathbf{A})$ is a category.
	
	$\text{Obj}(\text{Mor}(\mathbf{A})) = \text{Mor}\mathbf{A}$ (objects of category $\text{Mor}(\mathbf{A})$ are morphisms in $\mathbf{A}$).
	
	Let $\begin{aligned} & \quad \\ 
	& f:X \to Y \\
	& g: X' \to Y'\end{aligned}$, $f,g \in \text{Mor}(\mathbf{A})$ (i.e. $\begin{aligned} & \quad \\ 
	& f \in \text{Hom}(X,Y) \\ 
	&g \in \text{Hom}(X', Y') \end{aligned}$, for $X, Y, X', Y' \in \text{Obj}(\mathbf{A})$)
	
	Then
	\[
	\text{Hom}_{\text{Mor}(\mathbf{A})}(f,g) = \lbrace u :X \to X', v: Y \to Y'; g \circ u = v \circ f \rbrace
	\]
	Composition and identity in $\text{Mor}(\mathbf{A})$ are the obvious ones.
	
	So 
	\[
	\begin{gathered} 
	\text{Obj}(\text{Mor}(\mathbf{A})) = \text{Mor}(\mathbf{A}) \\
	\text{Mor}(\text{Mor}(\mathbf{A})) = \bigcup_{f,g \in \text{Mor}(\mathbf{A})} \text{Hom}(f,g) = \bigcup_{f,g \in \text{Mor}(\mathbf{A}) } \lbrace u : X\to X', v: Y \to Y' ; g \circ u = v \circ f \rbrace 
	\end{gathered} 
	\]
\[
	\begin{tikzpicture}
\matrix (m) [matrix of math nodes, row sep=5em, column sep=5em, minimum width=2em]
{
	X & Y \\
	X' & Y' \\
};
\path[->]
(m-1-1) edge node [above] {$f$} (m-1-2)
edge node [auto] {$u$} (m-2-1)
(m-1-2) edge node [auto] {$v$} (m-2-2)
(m-2-1) edge node [auto] {$g$} (m-2-2);
\end{tikzpicture}
\]
\end{definition}

cf. Def. 1.2.6 of Kashiwara and Schapira (2006) \cite{agat:KaSch2006}.

\begin{definition}
	\begin{enumerate}
		\item object $P \in \mathcal{C} \equiv \mathbf{A}$ is called \emph{initial} if $\forall \, X \in \mathbf{A}$, ($\equiv \forall \, x \in \text{Obj}(\mathbf{A})$), $\text{Hom}_{\mathbf{A}}(P, X) \simeq \lbrace \text{pt} \rbrace$. \\
		(Denote by $\emptyset_{\mathbf{A}}$ an \emph{initial object} in $\mathbf{A}$). \\
		(Note that if $P_1$ and $P_2$ are initial, then $\exists \, !$ isomorphism $P_1 \simeq P_2$)
		\item $P$ is terminal in $\mathbf{A}$ if $P$ is initial in $\mathbf{A}^{\text{op}}$, i.e. 
		
		$\forall \, X \in \mathbf{A}$, $\text{Hom}_{\mathbf{A}}(X, P) \simeq \lbrace \text{pt} \rbrace$. 
		
		Denote $\text{pt}_{\mathbf{A}}$ a terminal object in $\mathbf{A}$. 
		\item $P$ is zero ($0$) object if it's both initial and terminal. 
		
		Such a $P$ is denoted by $0$.
		
		If $\mathbf{A}$ has a zero object, $\forall \, $ object $X, Y \in \mathbf{A} \equiv \text{Obj}\mathbf{A}$, the morphism obtained as composition $X\to 0 \to Y$ is still denoted by $0: X\to Y$.
		
		(Note that composition of $0:X \to Y$, and any morphism $f:Y \to Z$ is $0:X \to Z$)
	\end{enumerate}
\end{definition}

cf. Example 1.2.7 of Kashiwara and Schapira (2006) \cite{agat:KaSch2006}. 
Example 
\begin{enumerate}
	\item[(i)] In category $\mathbf{Set}$, $\emptyset$ initial, $\lbrace \text{pt} \rbrace$ terminal. 
	\item[(ii)] Zero module $0$ is zero object in $\text{Mod}(R)$.
\end{enumerate}

Notation 1.2.8 of Kashiwara and Schapira (2006) \cite{agat:KaSch2006}:
\begin{enumerate}
	\item $\mathbf{Pt} \equiv $ category with a single object and a single morphism (the identity of this object)
	\item $\emptyset \equiv $ empty category with no objects (hence, no morphisms)
	\item $\bullet \to \bullet \equiv$ category which consists of 2 objects, say $a,b$, and 1 morphism, $a\to b$, other than $\text{id}_a, \text{id}_b \equiv 1_a, 1_b$. Denote this category by $\mathbf{Arr}$.  
\end{enumerate}

cf. Example 1.2.9 of Kashiwara and Schapira (2006) \cite{agat:KaSch2006}. 

Let $R$ be a ring. Let $N\in \text{Mod}(R^{\text{op}})$, $M\in \text{Mod}(R)$.

Category $\mathbf{C}$: 

$\text{Obj}\mathbf{C} \ni (f, L)$, where $L\in \text{Mod}(\mathbb{Z})$, $f$ bilinear map $f: N \times M \to L$ (i.e. it's $\mathbb{Z}$-bilinear and satisfies
\[
f(na, m) = f(n,am), \quad \, \forall \, a \in R
\])

Morphism from $f:N \times M \to L$ to $g: N\times M$ is a linear map $h:L\to K$ s.t. $h\circ f =g$. 

Since any bilinear map $f: N\times M \to L$ (i.e. any object of $\mathbf{C}$) factorizes uniquely through
\[
u  : N\times M \to N \otimes_R M
\]
object $(u, N \otimes_R M)$ is initial in $\mathbf{C}$

\section{Functors}

cf. Def. 1.2.10 of Kashiwara and Schapira (2006) \cite{agat:KaSch2006}, with terminology from pp. 13, Sec. 3. "Functors" from Mac Lane (1978) \cite{agat:MacL1978}

\begin{definition}[(covariant) Functor]\label{agat:Def:Functor}
\begin{enumerate}
	\item (covariant functor) Let categories $\mathbf{C}$, $\mathbf{D}$. \\
(covariant) \textbf{functor} $F: \mathbf{C} \to \mathbf{D}$ consists of (Mac Lane (1978) \cite{agat:MacL1978} says 2 suitably related functions)
\begin{itemize}
	\item map $F: \text{Obj}(\mathbf{C}) \to \text{Obj}(\mathbf{D})$ (i.e. $\forall \, C \in \text{Obj}(\mathbf{C})$, $F(C) \in \text{Obj}(\mathbf{D})$), \\
	(Mac Lane (1978) \cite{agat:MacL1978} calls this the \emph{object function} $T$ or $F$ in our notation; we'll call it the \emph{object map}) and 
	\item maps $F: \text{Hom}_{\mathbf{C}}(X,Y) \to \text{Hom}_{\mathbf{D}}(F(X), F(Y))$, so that 
		\[	
			F(f) : F(X) \to F(Y) \text{ or } F(f)(F(X)) = F(f(X)) 
		\] \, $\forall \, X,Y \in \text{Obj}(\mathbf{C})$ s.t.
	\begin{equation}
	\begin{gathered}
	F(1_X) = 1_{F(X)} \qquad \, \forall \, X \in \mathbf{C} \\
	F(g\circ f) = F(g) \circ F(f) \qquad \, \forall \, \begin{aligned} & \quad \\ 
	& f: X \to Y \\
	& g: Y \to Z \end{aligned}, \qquad \, X,Y,Z \in \text{Obj}(\mathbf{C})
	\end{gathered}
	\end{equation}
	Mac Lane (1978) \cite{agat:MacL1978} calls this the \emph{arrow function} (also written as $T$ for Mac Lane's notation; $F$ for our notation); we'll call $F$ to be the \emph{morphism map}. 
\end{itemize}
\item (composition law for functors) \\
For categories $\mathbf{A} , \mathbf{B}, \mathbf{C}$, functors $F: \mathbf{A} \to \mathbf{B}$, $G: \mathbf{B} \to \mathbf{C}$, \\
Composition $G\circ F : \mathbf{A} \to \mathbf{C}$, is a functor defined by 
\begin{equation}
\begin{aligned}
& (G\circ F)(X) = G(F(X)) \qquad \, \forall \, X \in \text{Obj}\mathbf{A}, \text{ and } \\
& (G\circ F)(f) = G(F(f)), \qquad \, \forall \, \text{ morphism } f \in \text{Mor}(\mathbf{C})
\end{aligned}
\end{equation}
\end{enumerate}
\end{definition}
Diagrammatically,

\[
\begin{gathered}
\begin{tikzpicture}
\matrix (m) [matrix of math nodes, row sep=1.8em, column sep=1.8em, minimum width=1.2em]
{
	X & Y \\
};
\path[->]
(m-1-1) edge node [above] {$f$} (m-1-2);
\end{tikzpicture} \quad \xrightarrow{F} \quad
\begin{tikzpicture}
\matrix (m) [matrix of math nodes, row sep=1.8em, column sep=1.8em, minimum width=1.2em]
{
	F(X) & F(Y) \\
};
\path[->]
(m-1-1) edge node [above] {$F(f)$} (m-1-2);
\end{tikzpicture} \\
\begin{tikzpicture}
\matrix (m) [matrix of math nodes, row sep=1.8em, column sep=1.8em, minimum width=1.2em]
{
	X & Y & Z \\
};
\path[->]
(m-1-1) edge node [above] {$f$} (m-1-2)
edge[bend right=45] node [below] {$g\circ f$} (m-1-3)
(m-1-2) edge node [above] {$g$} (m-1-3)
;
\end{tikzpicture} 
\quad \xrightarrow{F} \quad \begin{tikzpicture}
\matrix (m) [matrix of math nodes, row sep=1.8em, column sep=1.8em, minimum width=1.2em]
{
	F(X) & F(Y) & F(Z) \\
};
\path[->]
(m-1-1) edge node [above] {$F(f)$} (m-1-2)
edge[bend right=45] node [below] {$F(g\circ f)$} (m-1-3)
(m-1-2) edge node [above] {$F(g)$} (m-1-3)
;
\end{tikzpicture} 
\end{gathered}
\]
i.e.
\[
\begin{gathered} 
\begin{tikzpicture}
\matrix (m) [matrix of math nodes, row sep=3.8em, column sep=4.8em, minimum width=2.2em]
{
	X & Y \\
	F(X) & F(Y) \\
};
\path[->]
(m-1-1) edge node [above] {$f$} (m-1-2)
edge node [auto]  {$F$} (m-2-1)
(m-1-2) edge node [auto]  {$F$} (m-2-2)
(m-2-1) edge node [above] {$F(f)$} (m-2-2)        
;
\end{tikzpicture} \text{ and } 
\begin{tikzpicture}
\matrix (m) [matrix of math nodes, row sep=3.8em, column sep=4.8em, minimum width=2.2em]
{
	X & Y & Z \\
	F(X) & F(Y) & F(Z) \\ 
};
\path[->]
(m-1-1) edge node [above] {$f$} (m-1-2)
edge[bend left=45] node [above] {$g\circ f$} (m-1-3)
edge node [auto] {$F$} (m-2-1)
(m-1-2) edge node [above] {$g$} (m-1-3)
edge node [auto] {$F$} (m-2-2)
(m-1-3) edge node [auto] {$F$} (m-2-3)
(m-2-1) edge node [above] {$F(f)$} (m-2-2)
edge[bend right=45] node [below] {$F(g\circ f) = F(g) \circ F(f)$} (m-2-3)
(m-2-2) edge node [above] {$F(g)$} (m-2-3)  
;
\end{tikzpicture} 
\end{gathered}
\]

\subsection{Examples of Functors}

cf. pp. 13, 3. Functors, Maclane (1978) \cite{agat:MacL1978}

\emph{Power Set Functor}. 
\[
\begin{gathered}
	\begin{aligned}
		& P : \textbf{Set} \to \textbf{Set} \\
		& P:\text{Obj}\textbf{Set} \to \text{Obj}(\textbf{Set}) \\
		& P: X \in \text{Obj}\textbf{Set} \mapsto 2^X \in \text{Obj}\textbf{Set}, \text{ i.e. } \\
		& P:X \mapsto 2^X = \lbrace S | S \subseteq X \rbrace 
	\end{aligned} \\
P : \text{Hom}_{\textbf{Set}}(X,Y) \to \text{Hom}_{\textbf{Set}}(P(X), P(Y)) \\
\begin{aligned} 
& P(f): P(X) \to P(Y) \text{ or } P(f)(P(X)) = P(f(X)) \\
& P(f): S \in 2^X \mapsto f(S) \in 2^Y \text{ or } f(S) \subseteq Y 
\end{aligned} 
\end{gathered}
\]
\[
\begin{gathered}
P(1_X) = 1_{PX} \text{ and } P(g\circ f) = P(g) P(f), \text{ so $P$ defines a functor }
\end{gathered}
\]

Algebraic topology; singular homology in a given dim. $n$ ($n \in \mathbb{N}$) Maclane (1978) \cite{agat:MacL1978}

$\forall \, $ topological space $X$, assign $X \mapsto $ abelian group $H_n(X) \equiv n$th homology group of $X$, also \\
$\forall \, $ cont. $f:X \to Y$ of spaces corresponding group homomorphism. 
\[
\begin{gathered}
H_n(f) : H_n(X) \to H_n(Y) \\
\begin{aligned}
& H_n: \text{Obj}(\mathbf{Top}) \to \text{Obj}(\textbf{Ab}) \\
& H_n: X \mapsto H_n(X) 
\end{aligned} \\
H_n(f)(H_n(X)) = H_n(f(X))
\end{gathered}
\]
$\Longrightarrow H_n: \textbf{Top} \to \textbf{Ab}$.

e.g. if $X=Y = S^1$, $H_1(S^1) = \mathbb{Z}$

So group homomorphism $H_1(f): \mathbb{Z} \to \mathbb{Z}$ determined by integer $d\in \mathbb{Z}$ (image of 1) \\
$d$ is usual "degree" of cont. map $f:S^1 \to S^1$. 

In this case, and in general, homotopic maps $f,g : X \to Y$ \\
Then $\exists \,$ homotopy $H:X \times [0,1] \to Y$ s.t. $H(x,0) = f(x)$; $H(x,1) = g(x) \, \forall \, x \in X$




Example 1.2.3 (forgetful functors) Leinster (2014) \cite{agat:Lein2014}, pp. 18

\begin{enumerate}
	\item[(a)] $U : \textbf{Grp} \to \textbf{Set}$ \\
	if group $G$ then $U(G)$ is underlying set of $G$. \\
	If group homomorphism $f:G \to H$, \\
	$U(f)$ is function $f$ itself. \\
	So $U$ forgets group structure of groups and forgets group homomorphisms are homomorphisms.
	\item[(c)] $\textbf{Ab} \equiv $ category of abelian groups \\
	$\textbf{Ring} \to \textbf{Ab}$ forgets multiplicative structure. \\
	$U:\textbf{Ring} \to \textbf{Mon}$ forgets additive structure; $\textbf{Mon} \equiv $ category of monoids
\end{enumerate}




\begin{itemize}
	\item (monoid) homomorphism $h: M \to N$, $\forall \, $ monoid $M, N$ s.t. $\forall \, m,n \in M$, \\
	$h(m\cdot n) \equiv h(m\cdot_M n) = h(m) \cdot h(n) \equiv h(m) \cdot_N h(n)$ and \\
	$h(e_M) = e_N$. \\
	
	$h:m \mapsto h(m)$ so $h: \text{Obj}\mathbf{C}(M) \to \text{Obj}\mathbf{C}(N)$. \\
	Let $m\cdot n \equiv f_m(n)$, $f_m \in \text{Hom}\mathbf{C}(M)$
	\[
	\begin{aligned}
		& h(m) \cdot h(n) \equiv f_{h(m)}(h(n)), \, f_{h(m)} \in \text{Hom}\mathbf{C}(N) \\
 		& h(f_m) = h(m) \\
		& h(f_m)(h(n)) = h(m)\cdot h(n) = h(m\cdot n) = h(f_m\cdot n)
	\end{aligned}
	\]
	Compare the last statement with $F(f)(F(X)) = F(f(X))$.
	\item poset $P, Q$ (set $P$, equipped with $\leq$), poset as category $\mathbf{P}, \mathbf{Q}$, $\text{Obj}\mathbf{P} = P$, $\text{Obj}\mathbf{Q} = Q$ \\
	$\text{Hom}\mathbf{P} \ni m$ s.t. $m:x \mapsto y$ iff $x\leq y$. \\
	
	Let functor $F: \mathbf{P} \to \mathbf{Q}$ s.t. 
	\[
	\begin{aligned} 
		& F: \text{Obj}\mathbf{P} \to \text{Obj}\mathbf{Q} \\  
		&	F: x \mapsto F(x)
	\end{aligned} 
	\]
	\[
	\begin{aligned}
		& F: \text{Hom}\mathbf{P} \to \text{Hom}\mathbf{Q} \\ 
		& F: m \mapsto F(m)
	\end{aligned}
	\]
	s.t. for $m(x) =y$ and $x\leq y$, \\
	$F(m): F(x) \mapsto F(y)$ and $F(x) \leq F(y)$ (so that $F(m)$ is also monotone), i.e.
	\[
	F(m)(F(X)) = F(m(X)) = F(y)
	\]
\end{itemize}

Kleene closure itself is a functor from $\mathbf{Set}$ to $\textbf{Set}$, from $A$ to $A^*$, from $f$ to $f^*$. \\
It's the composition of $U\circ F$ of underlying functor $U:\textbf{Mon} \to \textbf{Set}$ and free functor $F:\textbf{Set} \to \textbf{Mon}$. 

\begin{itemize}
	\item \emph{arrow category} $\mathbf{C}^{\rightarrow}$ of category $\mathbf{C}$:
\[
\text{Obj}\mathbf{C}^{\rightarrow} = \text{Mor}\mathbf{C}
\]

Given $f:A \to B$, $f':A' \to B'$, morphism (arrow) $g: f \mapsto f'$ is a "commutative square", i.e. $g: \text{Hom}_{\mathbf{C}}(A\to B) \to \text{Hom}_{\mathbf{C}}(A' \to B')$

\[
\begin{gathered} 
\begin{tikzpicture}
\matrix (m) [matrix of math nodes, row sep=3.8em, column sep=4.8em, minimum width=2.2em]
{
	A & A' \\
	B & B' \\
};
\path[->]
(m-1-1) edge node [above] {$g_1$} (m-1-2)
edge node [auto]  {$f$} (m-2-1)
(m-1-2) edge node [auto]  {$f'$} (m-2-2)
(m-2-1) edge node [above] {$g_2$} (m-2-2)        
;
\end{tikzpicture} 
\end{gathered}
\]
where $g_1, g_2 \in \text{Mor}\mathbf{C}$, \\
i.e. such a morphism (arrow) is a pair of morphisms (arrows) $g=(g_1,g_2)$ in $\mathbf{C}$ s.t.
\[
g_2 \circ f = f'\circ g
\]

identity morphism (arrow) $1_f$ on object $f:A\to B$ is pair $(1_A, 1_B)$.

Composition of arrows is done componentwise:
\[
(h_1,h_2) \circ (g_1, g_2) = (h_1 \circ g_1, h_2\circ g_2)
\]

\[
\begin{tikzpicture}
\matrix (m) [matrix of math nodes, row sep=3.8em, column sep=4.8em, minimum width=2.2em]
{
	A & A' & A'' \\
	B & B' & B'' \\ 
};
\path[->]
(m-1-1) edge node [above] {$g_1$} (m-1-2)
edge[bend left=45] node [above] {$h_1\circ g_1$} (m-1-3)
edge node [auto] {$f$} (m-2-1)
(m-1-2) edge node [above] {$h_1$} (m-1-3)
edge node [auto] {$f'$} (m-2-2)
(m-1-3) edge node [auto] {$f''$} (m-2-3)
(m-2-1) edge node [above] {$g_2$} (m-2-2)
edge[bend right=45] node [below] {$h_2\circ g_2$} (m-2-3)
(m-2-2) edge node [above] {$h_2$} (m-2-3)  
;
\end{tikzpicture} 
\]

Observe that there are 2 functors:

\[
\begin{gathered} 
\begin{tikzpicture}
\matrix (m) [matrix of math nodes, row sep=3.8em, column sep=4.8em, minimum width=2.2em]
{
	\mathbf{C} & \mathbf{C}^{\rightarrow} & \mathbf{C} \\
};
\path[->]
(m-1-2) edge node [above] {$\textbf{dom}$} (m-1-1)
edge node [auto]  {$\textbf{cod}$} (m-1-3)
;
\end{tikzpicture} 
\end{gathered}
\]

\item \emph{algebra} cf. pp. 14, Sec. 3, "Functors" of Mac Lane (1978) \cite{agat:MacL1978}.

$\forall \, $ commutative ring $K$, set of all non-singular $n\times n$ matrices with entries in $K = $ general linear group $GL_n(K)$. \\
$\forall \, $ homomorphism $f: K \to K'$ of rings produces \\
homomorphism $GL_n f : GL_n(K) \to GL_n(K)$ of groups.  \\

$\forall \, n \in \mathbb{N}$ (natural numbers), defined functor $GL_n: \textbf{CRng} \to \textbf{Grp} \equiv GL_n \textbf{CommRing} \to \textbf{Grp}$.  \\


$\forall \, $ group $G$, set of all products of commutators $xyx^{-1} y^{-1}$ ($x,y \in G$) is a normal subgroup $[G, G]$ of $G$ called commutator subgroup. \\

Since $\forall \, $ homomorphism $G \to H$ of groups carries commutators to commutators.  \\
$G \mapsto [G, G]$ defines evident functor $\textbf{Grp} \to \textbf{Grp}$, \\
while $G \mapsto G/[G, G]$ define functor $\textbf{Grp} \to \textbf{Ab}$, factor-commutator functor. \\
TODO - understand the previous commutators.

\end{itemize} 

\subsubsection{Endofunctor}

\begin{definition}[Endofunctor]\label{agat:Def:Endofunctor}
Consider \textbf{endofunctor} $T : \mathbf{C} \to \mathbf{C}$. This implies that, for the \\

object map $\begin{aligned} & \quad \\ 
	& T : \text{Obj}{\mathbf{C}} \to \text{Obj}{\mathbf{C}} \\ 
	& T : X \to Y \end{aligned}$, \\
	
	object map $T$ acts like a morphism, in $\text{Mor}\mathbf{C}$ (!!!). In this case $T$ acts like $T\in \text{Hom}_{\mathbf{C}}(X, Y)$.  \\
	
	morphism map $T: \text{Hom}_{\mathbf{C}}(X, Y) \to \text{Hom}_{\mathbf{C}}(F(X), F(Y))$ s.t. 
	\begin{equation}
	\begin{gathered} 
		T(f) \equiv Tf:T(X) \to T(Y) \text{ so that } Tf \in \text{Hom}_{\mathbf{C}}(T(X), T(Y))  \\
		T(1_X) = 1_{T(X)} \in \text{Hom}_{\mathbf{C}}(T(X), T(Y)) (\text{implied from Kashiwara and Schapira (2006) \cite{agat:KaSch2006} in that} \, 1_X \in \text{Hom}(X,X)) \\
		T(g\circ f) = T(g) \circ T(f), \, \forall \, f:X \to Y, \, g:Y \to Z, \, X,Y,Z\in \text{Obj}{\mathbf{C}}
	\end{gathered} 
	\end{equation} 
	So $T(f) \equiv Tf$ is indeed a morphism map.

\end{definition}

\subsubsection{Hom functors}

Let $\mathbf{C}$ locally-small category (i.e. category s.t. hom-classes are actually sets and not proper classes).

$\forall \, A, B \in \text{Obj}\mathbf{C}$, 

\begin{definition}[covariant Hom functor]
	\textbf{covariant Hom functor} $\text{Hom}(A, -): \mathbf{C} \to \mathbf{Set}$
	\begin{equation}
	\begin{aligned}
	& \text{Hom}(A, - ) : \text{Obj}\mathbf{C} \to  2^{\text{Mor}(\mathbf{C})} \\ 
	& \text{Hom}(A, - ) : X \mapsto \text{Hom}(A,X) \\ 
	& \text{Hom}(A, - ) : \text{Mor}\mathbf{C} \to (\text{Mor}(\mathbf{C}) \to \text{Mor}(\mathbf{C})) \\
	& \text{Hom}(A, - ) : \text{Hom}(X,Y) \to (\text{Hom}(A,X) \to \text{Hom}(A,Y)) \\
	& \text{Hom}(A, - ) : f \mapsto \text{Hom}(A,f) \text{ where } \\
	& \begin{aligned} \text{Hom}(A,f) : & \text{Hom}(A,X) \to \text{Hom}(A,Y) \\ 
	& g\mapsto f\circ g \quad \, \forall \, g \in \text{Hom}(A,X) 
	\end{aligned} 
	\end{aligned}
	\end{equation}
\end{definition}

\begin{claim}
$\text{Hom}(A, -)$ is a functor.
\end{claim}

\begin{proof}
	\begin{itemize}
		\item \emph{identity}: $\text{Hom}(A,-) : 1_X \mapsto \text{Hom}(A, 1_X) \text{ where }$ 
		\[
		\begin{aligned}
			\text{Hom}(A, 1_X) : & \text{Hom}(A,X) \to \text{Hom}(A,X) \\ 
			& g\mapsto 1_X \cdot g = g , \quad \, \forall \, g\in \text{Hom}(A,X)
		\end{aligned}
		\]
		Hence
		\[
		\text{Hom}(A, -)(1_X) = 1_{\text{Hom}(A,X)} \in (\text{Hom}(A,X) \to \text{Hom}(A,X))
		\]
		\item \text{composition}: Let 
		\[
		\begin{aligned}
		& f \in \text{Hom}(X,Y) \\ 
		& g \in \text{Hom}(Y,Z) \\ 
		& k \in \text{Hom}(A,X) 
		\end{aligned}
		\]
		\[
		\text{Hom}(A, g) \circ \text{Hom}(A,f)(k) = \text{Hom}(A,g)(f\circ k) = g\circ f \circ k = (g\circ f)(k) = \text{Hom}(A, g\circ f)(k)
		\]
	\end{itemize}
\end{proof}

\begin{definition}[contravariant Hom functor]
	\textbf{contravariant Hom functor} $\text{Hom}(-, B): \mathbf{C} \to \mathbf{Set}$
	\begin{equation}
	\begin{aligned}
	& \text{Hom}(-, B ) : \text{Obj}\mathbf{C} \to  2^{\text{Mor}(\mathbf{C})} \\ 
	& \text{Hom}(-, B ) : X \mapsto \text{Hom}(X,A) \\ 
	& \text{Hom}(-, B ) : \text{Mor}\mathbf{C} \to (\text{Mor}(\mathbf{C}) \to \text{Mor}(\mathbf{C})) \\
	& \text{Hom}(-, B ) : \text{Hom}(X,Y) \to (\text{Hom}(Y,B) \to \text{Hom}(X,B)) \\
	& \text{Hom}(-, B ) : h \mapsto \text{Hom}(h,B) \text{ where } \\
	& \begin{aligned} \text{Hom}(h,B) : & \text{Hom}(Y,B) \to \text{Hom}(X,B) \\ 
	& g\mapsto g\circ h \quad \, \forall \, g \in \text{Hom}(Y,B) 
	\end{aligned} 
	\end{aligned}
	\end{equation}
\end{definition}

$\text{Hom}(A,-)$, $\text{Hom}(-,B)$ are related in a natural manner: for
\[
\begin{gathered}
\begin{aligned}
& f : B\to B', \, f\in \text{Hom}(B,B') \\ 
& g:A \to A', \, g\in \text{Hom}(A,A')
\end{aligned} \\
\begin{aligned}
& \text{Hom}(A,f) : \text{Hom}(A, B) \to \text{Hom}(A,B') \\ 
& \text{Hom}(g, B) : \text{Hom}(A,B) \to \text{Hom}(A', B)
\end{aligned}
\end{gathered}
\]

\[
\begin{gathered}
\begin{tikzpicture}
\matrix (m) [matrix of math nodes, row sep=4.8em, column sep=4.8em, minimum width=4.2em]
{
	\text{Hom}(A,B) & \text{Hom}(A',B) \\
	\text{Hom}(A,B') & \text{Hom}(A',B') \\
};
\path[->]
(m-1-1) edge node [above] {$\text{Hom}(g,B)$} (m-1-2)
 edge node [left] {$\text{Hom}(A,f)$} (m-2-1)
(m-1-2) edge node [right] {$\text{Hom}(A',f)$} (m-2-2)
(m-2-1) edge node [above] {$\text{Hom}(g,B')$} (m-2-2)
;
\end{tikzpicture}
\end{gathered}
\]

\begin{definition}[2-variable Hom functor]
	2-variable Hom functor $\text{Hom}(-,-)$ is a bifunctor.
	\[
	\text{Hom}(-,-) : \mathbf{C}^{\text{op}} \times \mathbf{C} \to \textbf{Set} 
	\]
	
\[
\begin{gathered}
\begin{tikzpicture}
\matrix (m) [matrix of math nodes, row sep=4.8em, column sep=4.8em, minimum width=4.2em]
{
	& D \\
	C & B \\
	A &  \\
};
\path[->]
(m-2-1) edge node [left] {$f\circ g \circ h$} (m-1-2)
edge node [above] {$ g \circ h$} (m-2-2)
edge node [left] {$  h$} (m-3-1)
(m-2-2) edge node [right] {$f$} (m-1-2)
(m-3-1) edge node [right] {$g$} (m-2-2)
;
\end{tikzpicture}
\end{gathered}
\]

$\text{Hom}(-,-)$ is a functor such that:
\[
\begin{gathered}
\begin{aligned} 
& \text{Hom}(h, f)(g) = f\circ g \circ h \\
& \text{Hom}(h, f) : \text{Hom}(A, B) \to \text{Hom}(C, D) 
\end{aligned} \\
\begin{aligned}
& h : C \to A, \, h \in \text{Hom}(C,A) \\ 
& f: B\dot D, \, f \in \text{Hom}(B,D)
\end{aligned} \\
(h,f) \in \text{Hom}(C,A) \times \text{Hom}(B,D) \mapsto \text{Hom}(h,f) \in (\text{Hom}(A,B) \to \text{Hom}(C,D)) \\
\text{Hom}(-,-) : \text{Hom}(C,A) \times \text{Hom}(B,D) \to (\text{Hom}(A,B) \to \text{Hom}(C,D) )
\end{gathered}
\]
and 
\[
\begin{aligned}
& \text{Hom}(-,-) : \text{Obj}\mathbf{C}^{\text{op}} \times \mathbf{C} \\ 
& \text{Hom}(-,-) : (C,D) \mapsto \text{Hom}(C,D)
\end{aligned}
\]
\end{definition}


\section{Construction of Categories}

\subsection{Slice category}

cf. Construction 4 of Awodey (2010) \cite{agat:Awod2010}, pp. 15.

\textbf{Slice category} $\mathbf{C}/C$, of category $\mathbf{C}$ over object $C \in \text{Obj}\mathbf{C}$, \\
$\text{Obj}\mathbf{C}/C \ni f \in \text{Mor}\mathbf{C}$ s.t. $\text{cod}(f) = C$ \\
$\text{Mor}\mathbf{C}/C \ni g$ from $f:X \to C$ to $f':X' \to C$ is a morphism (arrow) \, $g: X\to X'$ in $\mathbf{C}$ s.t. $f'\circ g = f$

\begin{equation}
\begin{gathered}
	\begin{tikzpicture}
\matrix (m) [matrix of math nodes, row sep=5em, column sep=5em, minimum width=2em]
{
	X & & X' \\
	& C &  \\
};
\path[->]
(m-1-1) edge node [above] {$g$} (m-1-3)
(m-1-1) edge node [below] {$f$} (m-2-2)
(m-1-3) edge node [below] {$f'$} (m-2-2);
\end{tikzpicture} \qquad \, 
	\begin{tikzpicture}
\matrix (m) [matrix of math nodes, row sep=5em, column sep=5em, minimum width=2em]
{
	X  \\
	C  \\
};
\path[->]
(m-1-1) edge node [left] {$f$} (m-2-2)
edge [loop right] node [right] {$1_{\mathbf{C}/C}$} (m-1-1);
\end{tikzpicture} \qquad \, 
	\begin{tikzpicture}
\matrix (m) [matrix of math nodes, row sep=5em, column sep=5em, minimum width=2em]
{
	X & X' & X'' \\
	& C &  \\
};
\path[->]
(m-1-1) edge node [above] {$g$} (m-1-2)
	edge node [left] {$f$} (m-2-1)
(m-1-2) edge node [above] {$h$} (m-1-3)
edge node [left] {$f'$} (m-2-2)
(m-1-3) edge node [right] {$f''$} (m-2-2);
\end{tikzpicture}
\end{gathered}
\end{equation}

If $\mathbf{C} = \mathbf{P}$ poset category, $p\in \mathbf{P}$ (i.e. $p \in \text{Obj}(\mathbf{P})$), then
\[
\mathbf{P} / p \simeq \downarrow (p)
\]
Slice category $\mathbf{P} / p$ is just the "principal ideal" $\downarrow(p)$ of elements $q\in \mathbf{P}$ with $q \leq p$.

\subsubsection{$S$-indexed set $X$ and indexed functions as morphisms of slice category $\mathbf{Set}/S$}

\begin{definition}[typed set]
	$S$-indexed set is set $X$ together with function $\tau : X \to S$. \\
	
	If $x\in X$, $\tau(x) = s$, then $x$ is of \textbf{type} $s$, so $X$ is a \textbf{typed set}. \\
	
	$\lbrace \tau^{-1}(s) | s \in S \rbrace$, $\tau^{-1}(s) \subset X \equiv $ family of sets indexed by $S$. \\
	
\end{definition} 
cf. 2.6.11 of Barr and Wells (2012) \cite{agat:BaWe2012}

\begin{definition}
	set $X$ typed by $S \to $ set $X'$ typed by $S$ that preserves typing (element of type $s\mapsto $ element of type $s$), is exactly an arrow (morphism) of slice category $\mathbf{Set}/S$, called \textbf{indexed function} or \textbf{typed function}.
\end{definition}

cf. 2.6.13 Indexed functions of Barr and Wells (2012) \cite{agat:BaWe2012}

\subsubsection{Underlying functor of the slice $F$}

Let $f: A\to C$, $f\in \text{Hom}(A,C) \subset \text{Mor}\mathbf{C}$ and $f\in \text{Obj}\mathbf{C}/C$.  \\

Let $f,g \in \text{Obj}\mathbf{C}/C$, s.t. $f\in \text{Hom}(A,C)$, $g\in \text{Hom}(B,C)$ \\
$h\in \text{Hom}(f,g) \subseteq \text{Mor}\mathbf{C}/C$ in that \\
$h : f\to g$ s.t. $h:A \to B$, so that $g\circ h = f$. \\

\textbf{underlying functor of the slice $F$} (notation $U$ in Barr and Wells (2012) \cite{agat:BaWe2012}):

\begin{equation}
\begin{aligned}
& F: \mathbf{C} /C \to \mathbf{C} \\ 
& F:\text{Obj}\mathbf{C}/C \to \text{Obj}\mathbf{C} \\ 
& F:f\mapsto F(f) = A \\ 
& F:\text{Mor}\mathbf{C}/C \to \text{Mor}\mathbf{C} \\ 
& F: h \mapsto h 
\end{aligned} 
\end{equation}

Special case: $\mathbf{C} = \textbf{Set}$, $C = S \ni \text{Obj}\textbf{Set}$, i.e. $S$ is a set.

For $f\in \text{Obj}\textbf{Set}/S$, $f: T \to S$, where $T$ is a set, object $f$ is an $S$-indexed set. (e.g. $x(i) \in T \subset \mathbb{R}$, where for some $y\in T \subset \mathbb{R}$, $\exists \, i \in \mathbb{Z}$ s.t. $y\mapsto i$).

\[
\begin{aligned}
& F:\textbf{Set}/S \to \textbf{Set} \\ 
&  F:f \mapsto T
\end{aligned}
\]

The underlying functor $F$ forgets the $S$-indexing.

\subsubsection{Coslice category}

\[
	\begin{tikzpicture}
\matrix (m) [matrix of math nodes, row sep=5em, column sep=5em, minimum width=2em]
{
	X & & X' \\
	& C &  \\
};
\path[->]
(m-1-3) edge node [above] {$\overline{g}$} (m-1-1)
(m-2-2) edge node [below] {$\overline{f}$} (m-1-1)
 edge node [below] {$\overline{f'}$} (m-1-3);
\end{tikzpicture} 
\]

Coslice category $C/\mathbf{C}$ of category $\mathbf{C}$ under object $C \in \text{Obj}\mathbf{C}$: \\
$\text{Obj}(C/\mathbf{C}) = \overline{f} \in \text{Mor}\mathbf{C}$ s.t. $\text{dom}\overline{f} = \overline{C} = C$, \\
$\text{Mor}(C/\mathbf{C}) \ni \begin{aligned} & \quad \\ & \overline{g} : \overline{X}' \to \overline{X} \text{s.t. } \\
& \overline{g} \circ \overline{f'} = \overline{f} \end{aligned}$ 

cf. Example 1.8 of Awodey (2010) \cite{agat:Awod2010}, pp. 15. \emph{pointed sets}

$\textbf{Sets}_*$ of pointed sets consisting of sets $A$ and distinguished element $a\in A$, and, i.e. \\
$\text{Obj}\textbf{Set}_* \ni $ set $A$ with distinguished element $a\in A$, \\
$\text{Mor}\textbf{Set}_* \ni f : (A, a) \to (B, b)$, function $f:A\to B$ preserves the "points" $f(a) = b$.

\[
\textbf{Sets}_* = \textbf{Set}_* \simeq 1 / \textbf{Sets}
\]
i.e. $\textbf{Sets}_*$ isomorphic to coslice category of sets "under" any singleton $1 = \lbrace * \rbrace$

\quad \\ 
functions $a: 1 \to A$ correspond uniquely to elements \\
\phantom{functions} $a(*) = a \in A$ \\
morphisms (arrows) $f:(A, a) \to (B, b)$ correspond exactly to 
\[
\begin{tikzpicture}
\matrix (m) [matrix of math nodes, row sep=5em, column sep=4em, minimum width=3em]
{
	1 & A \\ 
	&  B   \\
};
\path[->]
(m-1-1) edge node [above] {$a$} (m-1-2)
edge node [left] {$b$} (m-2-2)
(m-1-2) edge node [auto] {$f$} (m-2-2)
;
\end{tikzpicture} 
\]

\subsection{Free monoid}

Given set $S$, e.g. $S = $ "alphabet" $A$, \\
\emph{free monoid} $S^*$ is set $S^*$. \\

$S^* = $ set of all lists (finite sequences) of elements of $S$, \\
e.g. set of all \emph{words} (finite sequence of letters) \emph{over $A$ } $= A^*$. \\

\emph{Concatenation}: $* : \, \forall \, w , w^* \in A^*$, $w * w' = ww'$ \\
empty list (e.g. empty word "-"), $()$, $*$, or $\epsilon$ is a unit. Or we'll use this notation: \texttt{""}

\[
\Longrightarrow (S^*, *, \texttt{""}) \text{is a monoid. }
\]

monoid $M$ is \textbf{freely generated} by subset $A$ of $M$, if 
\begin{enumerate}
	\item "no junk": $\forall \, m \in M$, $m$ can be written as product of elements of $A$,
	\[
	m = a_1 \dots a_n, \, a_i \in A 
	\]
	\item "no noise": no "nontrivial" relations hold in $M$, i.e. if $a_1 \dots a_j = a_1' \dots a_k'$, then this is required by the axioms for monoid.
\end{enumerate}

$A^* = $ Kleene closure.

\subsubsection{Free monoid functor: underlying functor}

\[
\begin{aligned}
& F: \textbf{Set} \to \textbf{Mon} \\
& F: A \in \text{Obj}\textbf{Set} \mapsto \text{ free monoid } F(A) = \text{ Kleene closure } A^* 
\end{aligned}
\]
For $f:A \to B$, $f\in \text{Hom}_{\textbf{Set}}(A, B)$, 
\[
\begin{aligned}
& F: f\mapsto F(f) \text{ s.t. } \\
& F(f)(w) = F(f)(a_1 \dots a_m) = (f(a_1), \dots f(a_m)) = w' \in B^*, \text{ i.e. } \\
& F(f) = f^* : F(A) \to F(B)
\end{aligned}
\]


\section{Universal mapping property}

\subsection{Examples of Universal mapping property}

cf. Ch. 0 "Introduction" of Leinster (2014) \cite{agat:Lein2014}.

Example 0.1. Let $\lbrace 1 \rbrace$ denote set with 1 element. \\
Then $\forall \, $ sets $X \in \textbf{Set}$, $\exists \, ! $ map $X \to \lbrace 1 \rbrace$
\[
X \xrightarrow{ f} \lbrace 1 \rbrace 
\]

\begin{proof}
$\exists \, X \to \lbrace 1 \rbrace$, $\forall \, X$ because define $f: X \to \lbrace 1 \rbrace$ s.t. $f(x) = 1$, $\forall \, x \in X$. \\
(!) if for $X \to \lbrace 1 \rbrace$, then $x\mapsto 1$, so map is equal to $f$.
\end{proof} 

Example 0.2. $\forall \, $ rings $R$ with multiplicative identity $1$, $\exists \, !$ homomorphism $\mathbb{Z} \to R$. \label{agat:Example:integerringhomomorphism}

Define $\phi : \mathbb{Z} \to R$ by 
\[
\phi(n) = \begin{cases}  \sum_{i=1}^n 1 & \text{ if } n > 0 \\ 
0 & \text{ if } n = 0 \\
-\phi(-n) & \text{if } n < 0 \end{cases} 
\]

Check that $\phi$ homomorphism. (TODO)\\

(!) $\psi(1)= 1$ (homomorphisms preserve multiplicative identities) \\
homomorphisms preserve addition:
\[
\psi(n) = \psi(\sum_{i=1}^n 1) = \sum_{i=1}^n \psi(1) = \sum_{i = 1}^n 1 = \phi(n)
\]

homomorphisms preserve zero: $\psi(0) = 0 = \phi(0)$

homomorphisms preserve negatives: $\psi(n) = -\psi(-n) = -\sum_{i=1}^{-n} 1 = -\phi(-n) = \phi(n)$.

\begin{lemma}[0.3 Leinster (2014) \cite{agat:Lein2014}]
	Let ring $A$ s.t. $\forall \,$ rings $R, \exists \, !$ homomorphism $A \to R$ (i.e. $A$ is "initial"). Then $A \cong \mathbb{Z}$ (isomorphic). 	
\end{lemma}

\begin{proof}
		$A$ initial, so $\exists \, !$ homomorphism $\phi:A \to \mathbb{Z}$. \\
		$\mathbb{Z}$ initial from Ex. \ref{agat:Example:integerringhomomorphism}. $\exists \, !$ homomorphism $\phi' : \mathbb{Z} \to A$. \\
		
		$\phi' \circ \phi : A \to A$, $1_A : A \to A$ are homomorphisms. \\
		Since $A$ initial, $\phi' \circ \phi = 1_A$. \\
		$\phi \circ \phi' : \mathbb{Z}\to \mathbb{Z}$.  Since $\mathbb{Z}$ initial, $1_{\mathbb{Z}} : \mathbb{Z} \to \mathbb{Z} \Longrightarrow \phi \circ \phi' = 1_{\mathbb{Z}}$. 
		\[
		\Longrightarrow A \cong \mathbb{Z} \text{ (isomorphic) }
		\]
\end{proof}

Example 0.4. pp. 3 Leinster (2014) \cite{agat:Lein2014} Let vector space $V$ with basis $(v_s)_{s \in S}$ (.e.g if $V$ finite-dim., e.g. $S = \lbrace 1, 2 \dots , n \rbrace$).  \\
Thus $\forall \, $ vector space $W$, $\exists \, $ natural 1-to-1 correspondence between linear maps $\lbrace V \to W \rbrace$ and functions $S \to W$. \\
i.e. define $\begin{aligned} & \quad \\ 
& i : S \to V \\
& i(s) = v_s \quad \, (s \in S)\end{aligned}$ 

Then universal mapping property (given $i, f, \exists \, ! \overline{f}$)
\[
\begin{gathered} 
	\begin{tikzpicture}
\matrix (m) [matrix of math nodes, row sep=5em, column sep=4em, minimum width=3em]
{
	S & V  \\ 
	  & W   \\
};
\path[->]
(m-1-1) edge node [above] {$i$} (m-1-2)
 edge node [below] {$f$} (m-2-2)
;
\path[dashed,->]
(m-1-2)        edge node [right] {$\overline{f}$} (m-2-2)
;
\end{tikzpicture} \text{ or } \begin{tikzpicture}
\matrix (m) [matrix of math nodes, row sep=5em, column sep=4em, minimum width=3em]
{
	V & W  \\ 
	S &    \\
};
\path[dashed,->]
(m-1-1) edge node [above] {$\overline{f}$} (m-1-2)
;
\path[->]
(m-2-1) edge node [left] {$i$} (m-1-1)
edge node [below] {$f$} (m-1-2)
;
\end{tikzpicture}
\end{gathered} 
\]

\[
\begin{gathered}
\begin{aligned}
\lbrace S \to W \rbrace & \to \lbrace \text{ linear } V \to W \rbrace \\
\overline{f} \circ i & \mapsto \overline{f} 
\end{aligned} \quad \, \text{ bijective }
\end{gathered}
\]

Example 0.5. pp. 4 Leinster (2014) \cite{agat:Lein2014}. Given set $S$, topological space $D(S)$; equip $D(S)$ with discrete topology (all subsets open). \\
$\Longrightarrow \forall \, $ map $S \to $ space $X$ is cont. 

i.e. 
\[
\begin{gathered}
\begin{tikzpicture}
\matrix (m) [matrix of math nodes, row sep=5em, column sep=4em, minimum width=3em]
{
	D(S) & X  \\ 
	S &    \\
};
\path[dashed,->]
(m-1-1) edge node [above] {$\overline{f}$} (m-1-2)
;
\path[->]
(m-2-1) edge node [left] {$i$} (m-1-1)
edge node [below] {$f$} (m-1-2)
;
\end{tikzpicture}
\end{gathered}
\]

$\forall \, $ topological space $X$, $\forall \, f : S \to X$, $\exists \, !$ cont. $\overline{f} : D(S) \to X$ s.t. $\overline{f} \circ i = f$. \\
If $D(S)$ equipped with indiscrete topology (open $S = \lbrace \phi, S \rbrace$), then property is false.

Example 0.6. pp. 4 Leinster (2014) \cite{agat:Lein2014}. Given vector space $U, V, W$, bilinear map $f: U\times V \to W$ is linear. 
\[
\begin{gathered}
f(u, v_1 + \lambda v_2) = f(u, v_1) + \lambda f(u, v_2) \\
f(u_1 + \lambda u_2, v) = f(u_1, v) + \lambda f(u_2, v)
\end{gathered}
\]
$\forall \, u, u_1, u_2 \in U$, $v, v_1,v_2 \in V$, scalars $\lambda$. \\
e.g. scalar product (dot product) $\mathbf{u}\cdot \mathbf{v}$ bilniear, cross product bilinear. \\

$\exists \, $ "universal bilinear map out of $U\times V$", i.e. \\
$\exists \, $ vector space $T$, $\exists \, $ bilinear map $b:U\times V \to T$

\[
\begin{gathered}
\begin{tikzpicture}
\matrix (m) [matrix of math nodes, row sep=5em, column sep=4em, minimum width=3em]
{
	T & W  \\ 
	U \times V &    \\
};
\path[dashed,->]
(m-1-1) edge node [above] {$\overline{f}$} (m-1-2)
;
\path[->]
(m-2-1) edge node [left] {$b$} (m-1-1)
edge node [below] {$f$} (m-1-2)
;
\end{tikzpicture}
\end{gathered}
\]
i.e. $\forall \, $ bilinear $f$, given $b$, $\exists \, !$ linear $\overline{f}$, s.t. $\overline{f} \circ b = f$

\begin{lemma}[0.7, pp. 5 Leinster (2014) \cite{agat:Lein2014}]\label{agat:Lemma:UniveralMappingVectorSpaceProducts}
	Let vector spaces $U, V$. Suppose $\begin{aligned} & \quad \\ 
	& b : U \times V \to T \\
	& b' : U\times V \to T'\end{aligned}$ \quad , both universal bilinear maps. 
	
	Then $T \cong T'$, i.e. $\exists \, !$ isomorphism $j: T \to T'$ s.t. $j\circ b = b'$
\end{lemma}

\begin{proof}
	Take $b': U\times V \to T' \Longrightarrow \exists \, b$, $\exists \, !$ linear map $j: T \to T'$ s.t. $j\circ b = b'$. 
	
	Take $b:U \times V \to T$. $\exists \, b'$, $\exists \, !$ linear map $j':T' \to T$ s.t. $j'\circ b' = b$
	

\[
\begin{gathered}
\begin{tikzpicture}
\matrix (m) [matrix of math nodes, row sep=5em, column sep=4em, minimum width=3em]
{
	T & T'  \\ 
	U \times V & T   \\
};
\path[dashed,->]
(m-1-1) edge node [above] {$j$} (m-1-2)
(m-1-2) edge node [above] {$j'$} (m-2-2)
;
\path[->]
(m-2-1) edge node [left] {$b$} (m-1-1)
edge node [below] {$b'$} (m-1-2)
edge node [below] {$b$} (m-2-2)
;
\end{tikzpicture}
\end{gathered}
\]
Now linear $j'\circ j : T\to T$, s.t. $(j'\circ j) \circ b = b$. 

But also $1_T : T\to T$ linear and $1_T \circ b = b$

\[
\Longrightarrow j' \circ j = 1_T
\]
i.e. 
\[
\begin{gathered}
\begin{tikzpicture}
\matrix (m) [matrix of math nodes, row sep=5em, column sep=4em, minimum width=3em]
{
	T & T  \\ 
	U \times V &    \\
};
\path[dashed,->]
(m-1-1) edge node [above] {$j'\circ j$} (m-1-2)
;
\path[->]
(m-2-1) edge node [left] {$b$} (m-1-1)
edge node [below] {$b$} (m-1-2)
;
\end{tikzpicture}
\qquad \, \text{ but so is } \begin{tikzpicture}
\matrix (m) [matrix of math nodes, row sep=5em, column sep=4em, minimum width=3em]
{
	T & T  \\ 
	U \times V &    \\
};
\path[dashed,->]
(m-1-1) edge node [above] {$1_T$} (m-1-2)
;
\path[->]
(m-2-1) edge node [left] {$b$} (m-1-1)
edge node [below] {$b$} (m-1-2)
;
\end{tikzpicture}
\end{gathered}
\]

Similarly for $j\circ j' = 1_{T'}$. 

So $j$ is an isomorphism.

\end{proof} 

By Lemma \ref{agat:Lemma:UniveralMappingVectorSpaceProducts}, $\exists \, !$ tensor product $U\otimes V$, not "a" tensor product.

Example 0.8 of Leinster (2014) \cite{agat:Lein2014}. Let $\theta : G \to H$ be homomorphism of groups.

\[
\begin{tikzcd}
\text{ker}{(\theta)} \arrow[r, hook, "i"] & G \arrow[r, shift left, "\theta"] \arrow[r, shift right, "\epsilon"]& H \\
\end{tikzcd}
\]

where $i$ inclusion, $\epsilon$ trivial homomorphism, i.e.$i(x) = x \, \forall \, x \in \text{ker}(\theta)$, $\epsilon(g) = 1$ \, $\forall \, g \in G$.

map $i$ into $G$ s.t. $\theta \circ i = \epsilon \circ i$ and is universal (Ex. 0.11)

Example 0.9 of Leinster (2014) \cite{agat:Lein2014}. Let topological space $X = U \cup V$ be covered by 2 open subsets $U, V$.

\[
\begin{tikzcd}[sep=large] 
	U \cap V \arrow[r, hook, "i"] \arrow[d, hook, "j"] & U \arrow[d, hook, "j'"] \\
	V \arrow[r, hook, "i'"] & X
\end{tikzcd}
\]
has universal property
\[
\begin{tikzcd}[sep=large] 
U \cap V \arrow[r, hook, "i"] \arrow[d, hook, "j"] & U \arrow[d, hook, "j'"] \arrow[ddr, "f"] & \\
V \arrow[r, hook, "i'"] \arrow[drr, "g"] & X \arrow[dr, dotted, "h"] & \\
& & Y 
\end{tikzcd}
\]
$\forall \, g, Y, f$ s.t. $f\circ i = g\circ j$, $\exists ! \, $ cont. $h: X\to Y$ s.t. $h\circ j' = f$, $h\circ i' = g$.

Under favorable conditions, induced diagram of fundamental groups.

\[
\begin{tikzcd}[sep=large] 
\pi_1(U \cup V) \arrow[r, "i_*"] \arrow[d, "j_*"] & \pi_1(U) \arrow[d, "j'_*"] \\
\pi_1(V) \arrow[r, "i'_*"] & \pi_1(X)
\end{tikzcd}
\]
has same universal property: Van Kampen's thm.

\[
\begin{aligned}
& \theta : G \to H \\ 
& \epsilon : G \to H
\end{aligned}
\]

$\theta, \epsilon$ have common domain and codomain.

$i: \text{ker}(\theta) \to G$ is s.t. $\theta \circ i = \epsilon \circ i$

\[
\begin{tikzcd}
\text{ker}{(\theta)} \arrow[r, hook, "i"] & G \arrow[r, shift left, "\theta"] \arrow[r, shift right, "\epsilon"]& H \\
J \arrow[u, dotted, "u"] \arrow[ur, "a"] & & \\
\end{tikzcd}
\]
$\forall \, a : J \to G$ s.t. $\theta \circ a = \epsilon \circ a$, then $\exists \, ! \, u : J \to \text{ker}{\theta}$ s.t. 
\[
a = i \circ u 
\]

\subsubsection{Free monoid $A^*$ on set $A$}

$\exists \, i : A \to A^*$, given any monoid $N$, any $f: A \to |N|$ (where $|N|$ is the underlying set of monoid $N$), \\
$\exists \, ! $ monoid homomorphism $\overline{f}: A^* \to N$ s.t. $|\overline{f}| \circ i = f$.


\[
\begin{aligned}
\textbf{Mon} : & 
\begin{tikzpicture}
\matrix (m) [matrix of math nodes, row sep=5em, column sep=5em, minimum width=2em]
{
	A^* &  N \\
};
\path[->]
(m-1-1) edge [dotted] node [above] {$\overline{f}$} (m-1-2);
\end{tikzpicture}  \\
\quad \\ 
\textbf{Sets} : & \begin{tikzpicture}
\matrix (m) [matrix of math nodes, row sep=5em, column sep=5em, minimum width=2em]
{
	\lvert A^* \rvert &  \lvert N \rvert \\
	A & \\
};
\path[->]
(m-1-1) edge [dotted] node [above] {$|\overline{f}|$} (m-1-2)
(m-2-1) edge node [left] {$i$} (m-1-1)
edge node [right] {$f$} (m-1-2);
\end{tikzpicture}  
\end{aligned} 
\]

\begin{proposition}[Proposition 1.9, Awoday (2010) \cite{agat:Awod2010}]] $A^*$ has universal mapping property of free monoid on $A$.
	\end{proposition} 
\begin{proof}
	Givne $f: A \to |N|$, define $\overline{f} : A^* \to N$, by 
	\[
	\begin{aligned}
	& \overline{f}("") = e_N \\ 
	& \overline{f}(a_1 \dots a_n) = f(a_1) \dots f(a_n) 
	\end{aligned}
	\]	
	$\overline{f}$ is thus a homomorphism with $\overline{f}(a) = f(a)$ \quad \, $\forall \, a \in A$. \\
	
	If $g: A^* \to N$ s.t. $g(a) =f(a)$ \quad \, $\forall \, a \in A$, then $\forall \, (a_1, \dots a_n) \in A^*$, 
	\[
	\overline{f}(a_1 \dots a_n) = \overline{f}(a_1) \dots \overline{f}(a_n) = f(a_1) \dots f(a_n) = g(a_1) \dots g(a_n) = g(a_1 \dots a_n)
	\]
	since $\overline{f} = g$, $\overline{f}$ unique.
\end{proof} 

Existence part of UMP is "no noise." \\
Unique part of UMP is "no junk."

TODO Prop. 1.10, Awodey.






\exercisehead{1} cf. 3.1.23 Exercises of Barr and Wells (2012) \cite{agat:BaWe2012}

Given semigroup $S$, construct monoid $M = S \cup \lbrace e \rbrace$, $e\notin S$. \\
Let $e=\lbrace S \rbrace$. \\
Multiplication in $M$ defined as 
\begin{enumerate}
	\item $xy \in S$ if $x,y \in S$ 
	\item $xe=ex = x \quad \, \forall \, x \in M$
\end{enumerate}
$M \equiv S^1$ in semigroup literature.

\begin{enumerate}
	\item $S^1$ is a monoid (Show). \\
	If $a,b,c \in S \subset M$, then $(ab)c = a(bc)$ (associative holds by definition of semigroup). \\
	$\forall \, a \in S$, $ae = ea =a$, by how we defined multiplication. \\
	$ee = ee =e $ by either definition, or that multiplication in semigroup $S$ is closed.
	\[
	\lbrace S \rbrace \lbrace S \rbrace = \lbrace S \rbrace \lbrace S \rbrace = \lbrace S \rbrace
	\]
	
	If for $a,b,c \in M$, any of $a,b$ or $c=e$, use multiplication to show equality for $(ab)c = a(bc)$; e.g.
	\[
	\begin{aligned}
	& (eb) c = (b)\cdot c = bc = e(bc) = bc \\ 
	& (ae) c = a\cdot c =a \cdot (ec) = a(ec)
	\end{aligned}
	\]
	$S^1$ is a monoid.
	\item 
\end{enumerate}

\section{Actions, Finite State Machines}

\begin{definition}
Let monoid $M$ with identity $1$, and a set $S$. \\
\textbf{action} of $M$ on $S$, $\alpha$, \\
\begin{equation}
\begin{aligned}
& \alpha : M \times S \to S \text{ s.t. } \\  
& \alpha(1,s) = s \quad \, \forall \, s \in S \\
& \alpha(mn, s) = \alpha(m, \alpha(n, s)) \quad \, \forall \, m,n \in M, \, s\in S
\end{aligned}
\end{equation}

Write $ms \equiv \alpha(m, s)$ then 
\[
\begin{gathered} 
1s = s \\
(mn) s = m (ns) \quad \, \forall \, m, n \in M, \, s \in S
\end{gathered}
\]
\end{definition}

\begin{definition}[equivariant map]
	Let monoid $M$ with actions on sets $S,T$. \\
	\textbf{equivariant map} $\phi:S \to T$ s.t. $m\phi(s) = \phi(ms)$
\end{definition}

Suppose 2 equivariant maps $\phi_{TS}$, $\phi_{UT}$, sets $S,T,U$ s.t. monoid $M$ has actions on $S,T,U$: 
\[
\begin{aligned}
& \phi_{TS} : S \to T \\ 
& \phi_{UT} : T \to U
\end{aligned}
\]

\[
m\phi_{UT} \circ \phi_{TS}(s) = m \phi_{UT}(\phi_{TS}(s)) = \phi_{UT}(m\phi_{TS}(s)) = \phi_{UT} \phi_{TS} (ms) 
\]

Let $\phi_{UT} \circ \phi_{TS} \equiv \phi_{US}$. $\phi_{US}$ equivariant.

Associativity: given equivariant maps 
\[
\begin{aligned}
& \phi_{TS} : S \to T \\ 
& \phi_{UT} : T \to U \\ 
& \phi_{VU} : U \to V 
\end{aligned}
\]

\[
m \phi_{VU} \phi_{UT} \phi_{TS}(s) = m \phi_{VU} (\phi_{UT} \phi_{TS}(s)) = \phi_{VU} (m \phi_{UT} \phi_{TS}(s)) = (\phi_{VU} \phi_{UT}) (\phi_{TS}(ms))
\]

Thus, equivariant maps obey associativity on $\forall \, S \in \text{Obj}(M-\mathbf{Act})$, (set $S$ that $M$ has action on) \\
category $M-\mathbf{Act}$, $\text{Obj}(M-\mathbf{Act}) = \lbrace \text{ sets that monoid $M$ has action on }\rbrace$ \\
\phantom{category $M-\mathbf{Act}$,} $\text{Mor}(M-\mathbf{Act}) = \lbrace \text{ equivariant maps $\phi$ s.t. } \forall \, m \in M, \, m\phi(s) = \phi(ms) \rbrace$.

\subsection{Actions as functors}

Let $\alpha$ action of monoid $M$ on set $S$. \\

Consider $\mathbf{C}(M)$ (monoid as a category; i.e. $\text{Obj}\mathbf{C}(M) = \lbrace M \rbrace, \, \text{Mor}\mathbf{C}(M) = M$). \\
action $\alpha$ determines function $F_{\alpha} : \mathbf{C}(M) \to \mathbf{Set}$, 
\begin{enumerate}
	\item $F_{\alpha}(*) = S \equiv F_{\alpha}(\lbrace M \rbrace ) = S$
	\item  $F_{\alpha}(m) = s \mapsto \alpha(m, s) \, \forall \, m \in M, s \in S$
\end{enumerate}
i.e.

\[
\begin{gathered}
\begin{aligned}
& F_{\alpha} : \text{Obj}\mathbf{C}(M) \to \text{Obj}\mathbf{Set} \\ 
& F_{\alpha} : \lbrace M \rbrace \mapsto S 
\end{aligned} \\
\begin{aligned}
& F_{\alpha} : \text{Mor}\mathbf{C}(M) \to \text{Mor}\textbf{Set} \\ 
& F_{\alpha} : m \mapsto (s\mapsto \alpha(m,s)) \in \text{Hom}(S,S)
\end{aligned}
\end{gathered}
\]

Barr and Wells (1998) \cite{agat:BW1998} uses this notation: $\mathcal{M} = (A, S, s_0, \phi)$ for a "machine." Compare this to the notation used in these notes:
\[
\text{FSM} = (\Sigma, S, s_0, \delta)
\]

Take note that the finite set or input alphabet $A \equiv \Sigma$ is \textbf{not} a monoid. \\

$\forall \, $ string $\equiv $ finite sequence $\equiv $ "word" induces sequence of transitions in FSM starting at $s_0$, and ending on some final state; precisely, define
\[
\begin{aligned}
& \delta^* : \Sigma^* \times S \to S \\ 
& \delta^*("", s) = s \quad \, \forall \in s \in S \\ 
& \delta^*((a)w, s) = \delta(a, \delta^*(w,s)) \quad \, \forall \, s \in S, \, w\in \Sigma^*, \, a\in A
\end{aligned}
\]

Recall free monoid $\Sigma^*$, $(\Sigma^*, *, "")$ where $\Sigma^* \in \text{Obj}\mathbf{Set}$ is a set. \\

\begin{proposition}
	$\delta^*$ is an action of $\Sigma^*$ on $S$.
\end{proposition}

\begin{proof}
	Assume
	\[
	\delta^*(wv, s) = \delta^*(w, \delta^*(v,s))
	\]	
	then
	\[
	\begin{gathered}
		\delta^*((a)wv, s) = \delta(a, \delta^*(wv, s)) = \delta(a, \delta^*(w, \delta^*(v,s))) = \delta^*(aw, \delta^*(v,s)) 
	\end{gathered}
	\]
1st., 3rd. equalities from $\delta^*((a)w, s) = \delta(a, \delta^*(w,s))$ (definition of $\delta^*$). \\
2nd. equality is from inductive hypothesis.
\end{proof}

\textbf{recognizer} - subset $L\subseteq A^*$ of strings which drive FSM from $s_0$ (start state) to acceptor state, is then set of strings, or language, which is \textbf{recognized} by machine FSM, this is the machine as \textbf{recognizer}.

\textbf{transducer} - FSM outputs string of symbols (not necessarily in same alphabet) $\forall \, $ state it enters, or each transition it undergoes.

\subsection{Set-valued functors as actions}

cf. 3.2.6 of Barr and Wells (2012) \cite{agat:BaWe2012}. 


\section{Products, Coproducts}

\subsection{Sources}

It appears Ad\'{a}mek, Herrlich, and Strecker (2004) \cite{agat:AHS2004} defines \emph{sources} to simply give a name and formalize a tuple.  

\begin{definition}[source]
	\textbf{source} is a tuple: $(a, (f_i)_{i\in I})$, $f_i:A\to A_i$
\end{definition}

\subsection{Products}

\begin{definition}[Products]\label{agat:Def:Products}
	
	(in Turi's notation \cite{agat:Turi2001})
	
	Given objects $C_1,C_2$ of category $\mathbb{C}$, \textbf{product} (if exists) consists of object $C_1 \times C_2$ of $\mathbb{C}$ and $\begin{aligned} & \quad \\
	& \pi_1 : C_1 \times C_2 \to C_1 \\
	& \pi_2: C_1 \times C_2 \to C_2 \end{aligned}$ s.t. \\
	$\forall \, $ object $A$ of $\mathbb{C}$, \, $\forall \, \begin{aligned} & \quad \\
	& f : A \to C_1 \\
	& g : A \to C_2 \end{aligned}$ \quad \, $\exists \, ! \quad \, \langle f,g \rangle : A \to C_1 \times C_2$ s.t. $\begin{aligned} & \quad \\
	& f = \pi_1 \circ \langle f,g \rangle \\
	& g = \pi_2 \circ \langle f,g \rangle \end{aligned}$, i.e. 
	
	\begin{tikzpicture}
	\matrix (m) [matrix of math nodes, row sep=5em, column sep=4em, minimum width=3em]
	{
		& A & \\ 
		C_1  & C_1 \times C_2  &  C_2  \\
	};
	\path[->]
	(m-1-2) edge node [above] {$f$} (m-2-1)
	edge node [above] {$g$} (m-2-3)
	(m-2-2) edge node [below] {$\pi_1$} (m-2-1)
	edge node [below] {$\pi_2$} (m-2-3)
	;
	\path[dashed] 
	(m-1-2) edge node [right] {$\langle f,g\rangle$} (m-2-2)
	;
	\end{tikzpicture} 
	
	(compare with Leinster (2014) \cite{agat:Lein2014})
	
	Let category $\mathcal{A}$, $X,Y \in \mathcal{A}$, \textbf{product} of $X,Y$ consists of object $P$ and maps
	
	(compare this definition with Ad\'{a}mek, Herrlich, and Strecker (2004) \cite{agat:AHS2004} and their notation)
	
	\textbf{product} consisting of 
	\[
	\begin{aligned}
	C_1 \times C_2 \times \dots \times C_{\mathcal{N}} \in \text{Obj}\mathbf{C} \\
	\pi_1 : C_1 \times C_2 \times \dots \times C_{\mathcal{N}} \to C_1 \\ 
	\pi_2 : C_1 \times C_2 \times \dots \times C_{\mathcal{N}} \to C_2 \\ 
	\vdots \\
	\pi_{\mathcal{N}} : C_1 \times C_2 \times \dots \times C_{\mathcal{N}} \to C_{\mathcal{N}} \\ 
	\end{aligned}
	\]
	
	is s.t. 
	
	$\forall \, \begin{aligned}  \quad \\ 
	A \in \text{Obj}\mathbf{C} \\
	f_1 : A \to C_1 \\
	f_2 : A \to C_2 \\
	\vdots \\
	f_{\mathcal{N}} : A \to C_{\mathcal{N}} \end{aligned}$, 
	
	$\exists \, ! \langle f_1 ,f_2 , \dots , f_{\mathcal{N}} \rangle : A \to C_1 \times C_2 \times \dots \times C_{\mathcal{N}}$ s.t.
	\[
	\begin{aligned}
	f_1 = \pi_1 \circ \langle f_1 , f_2 , \dots f_{\mathcal{N}} \rangle \\ 
	f_2 = \pi_2 \circ \langle f_1 , f_2 , \dots f_{\mathcal{N}} \rangle \\ 
	\vdots  \\ 
	f_{\mathcal{N}} = \pi_{\mathcal{N}} \circ \langle f_1 , f_2 , \dots f_{\mathcal{N}} \rangle 
	\end{aligned}
	\]
\end{definition}

Let's use the notation of Ad\'{a}mek, Herrlich, and Strecker (2004) \cite{agat:AHS2004}. Also note that these references \emph{assume} the \emph{universal mapping property} in the definition of a \textbf{product}.

\begin{definition}[Product (assuming univeral mapping property)]
	\begin{equation}
\begin{gathered}
\begin{tikzpicture}
\matrix (m) [matrix of math nodes, row sep=4em, column sep=4em, minimum width=3em]
{
	& & A \\
	P & A_i & \\
	& A_{i'} & \\
};
\path[->]
(m-2-1) edge node [above] {$p_i$} (m-2-2)
edge node [above] {$p_{i'}$} (m-3-2)
(m-1-3) edge node [above] {$f_i$} (m-2-2)
edge node [above] {$f_{i'}$} (m-3-2)
;
\path[dashed,->]
(m-1-3) edge node [above] {$f$} (m-2-1)
;
\end{tikzpicture} 
\end{gathered}
\end{equation}
For the case of only 2 objects associated with $I$, 
	\begin{equation}
\begin{gathered}
\begin{tikzpicture}
\matrix (m) [matrix of math nodes, row sep=5em, column sep=4em, minimum width=3em]
{
	& A  & \\ 
	A_1  & P  & A_2   \\
};
\path[->]
(m-1-2) edge node [above] {$f_1$} (m-2-1)
edge node [above] {$f_2$} (m-2-3)
(m-2-2) edge node [above] {$p_1$} (m-2-1)
edge node [above] {$p_2$} (m-2-3)
;
\path[dashed,->]
(m-1-2) edge node [right] {$f$} (m-2-2)
;
\end{tikzpicture} 
\end{gathered}
\end{equation}


\textbf{product} $P, \lbrace p_i, A_i \rbrace_{i\in I}$ s.t. $\forall \, A, \lbrace f_i, A_i \rbrace_{i \in I}$, $\exists \, !$ morphism $f:A \to P$ s.t. $f_i = p_i \circ f$, $i\in I$.
\end{definition}
This definition is the same one used in pp. 35, Def. 2.16, of Awodey (2010) \cite{agat:Awod2010}.

Consider the definition that doesn't use the universal mapping property from Pareigis (2004) \cite{agat:Pare2004}. 

\begin{definition}[Product without universal property]
	Given objects $A, B \in \text{Obj}\mathbf{C}$, category $\mathbf{C}$. 
	
	Object $A\times B$, morphisms $p_A, p_B$, $\begin{aligned} & \quad \\ 
		& p_A : A \times B \to A \\ 
		& p_B : A\times B \to B
		\end{aligned}$ is called a (not the) (categorical) product of $A, B$, projections $p_A, p_B$. 
		
		If $\forall \, $ object $T\in \text{Obj}\mathbf{C}$, $\exists \, $ isomorphism

		\begin{equation}
			\text{Mor}_{\mathbf{C}}(T, A \times B)  \cong \text{Mor}_{\mathbf{C}}(T, A) \times \text{Mor}_{\mathbf{C}}(T,B)		
		\end{equation}
(where $\text{Mor}_{\mathbf{C}}$ is the Cartesian product) s.t.

\begin{equation}\label{agat:Eq:ProductDefinitionWithoutUniversalPropertyDiagram}
\begin{gathered}
\begin{tikzpicture}
\matrix (m) [matrix of math nodes, row sep=5em, column sep=5em, minimum width=4em]
{
	& \text{Mor}_{\mathbf{C}}(T, A\times B)  & \\ 
	\text{Mor}_{\mathbf{C}}(T, A)  & \text{Mor}_{\mathbf{C}}(T, A) \times \text{Mor}_{\mathbf{C}}(T, B)  & \text{Mor}_{\mathbf{C}}(T, B)   \\
};
\path[->]
(m-1-2) edge node [left] {$\text{Mor}_{\mathbf{C}}(T,p_A)$} (m-2-1)
edge node [right] {$\text{Mor}_{\mathbf{C}}(T,p_B)$} (m-2-3)
(m-2-2) edge node [above] {$p_{\text{Mor}_{\mathbf{C}}(T, A)}$} (m-2-1)
edge node [above] {$p_{\text{Mor}_{\mathbf{C}}(T, B)}$} (m-2-3)
;
\path[dashed,->]
(m-1-2) edge node [right] {$h$} (m-2-2)
;
\end{tikzpicture} \\
\quad \\ 
\begin{tikzpicture}
\matrix (m) [matrix of math nodes, row sep=5em, column sep=5em, minimum width=4em]
{
	& u  & \\ 
	f  & h(u) = (f,g)  & g   \\
};
\path[|->]
(m-1-2) edge node [left] {$\text{Mor}_{\mathbf{C}}(T,p_A)$} (m-2-1)
edge node [right] {$\text{Mor}_{\mathbf{C}}(T,p_B)$} (m-2-3)
(m-2-2) edge node [above] {$p_{\text{Mor}_{\mathbf{C}}(T, A)}$} (m-2-1)
edge node [above] {$p_{\text{Mor}_{\mathbf{C}}(T, B)}$} (m-2-3)
;
\path[dashed,|->]
(m-1-2) edge node [right] {$h$} (m-2-2)
;
\end{tikzpicture} 
\end{gathered}
\end{equation}
where $\begin{aligned} & \quad \\ 
& \text{Mor}_{\mathbf{C}}(T, p_A)(u) = f = p_{\text{Mor}_{\mathbf{C}}(T, A)} h(u) \\
& \text{Mor}_{\mathbf{C}}(T, p_B)(u) = g = p_{\text{Mor}_{\mathbf{C}}(T, B)} h(u) \end{aligned}$
\end{definition}

\begin{proposition}[Equivalent definition of product by universal mapping property]
Given objects $A, B \in \text{Obj}\mathbf{C}$, category $\mathbf{C}$, object $A\times B$, morphisms $\begin{aligned} & \quad \\ 
	p_A: A\times B \to A \\
	p_B: A\times B \to B \end{aligned}$ is a (categorical) product, \\
	
iff $\forall \, $ object $T \in \text{Obj}\mathbf{C}$, $\forall \, \begin{aligned} & \quad \\ 
& f: T\to A \\
& g: T \to B \end{aligned}$, $\exists \, !$ morphism $(f,g) : T \to A\times B$, s.t. 

\begin{equation}
\begin{tikzpicture}
\matrix (m) [matrix of math nodes, row sep=5em, column sep=4em, minimum width=3em]
{
	& T  & \\ 
	A  & A\times B  & B   \\
};
\path[->]
(m-1-2) edge node [left] {$f$} (m-2-1)
edge node [right] {$g$} (m-2-3)
(m-2-2) edge node [above] {$p_A$} (m-2-1)
edge node [above] {$p_B$} (m-2-3)
;
\path[dashed,->]
(m-1-2) edge node [left] {$(f,g)$} (m-2-2)
;
\end{tikzpicture} 
\end{equation}	
\end{proposition}
cf. Proposition "Characterization of products by universal mapping property", 2.7.4 in Pareigis (2004) \cite{agat:Pare2004}.

\begin{proof}
$\Longleftarrow$ : \\

	Let $(A\times B, p_A, p_B)$ be a "product" of $A\times B$, obeying the universal (mapping) property. Let $T\in \text{Obj}\mathbf{C}$. Let $u\in \text{Mor}_{\mathbf{C}}(T, A\times B)$. \\
	
	Define $\begin{aligned} & \quad  \\
	& f:= p_A \circ u : T \to A \\ 
	& g:= p_b \circ u : T \to B\end{aligned}$  \\
	Then $(f,g) \in \text{Mor}_{\mathbf{C}}(T,A) \times \text{Mor}_{\mathbf{C}}(T,B)$. \\
	
	Define $h(u):= (f,g) \in \text{Mor}_{\mathbf{C}}(T,A) \times \text{Mor}_{\mathbf{C}}(T,B)$. \\
	Then, 
	\[
	\begin{aligned}
		& \text{Mor}_{\mathbf{C}}(T,p_A)(u) = p_A \circ u = f = p_{\text{Mor}_{\mathbf{C}}(T,A)} \circ h(u) \\	
		& \text{Mor}_{\mathbf{C}}(T,p_B)(u) = p_B \circ u = g = p_{\text{Mor}_{\mathbf{C}}(T,B)} \circ h(u) \\	
	\end{aligned}	
	\]
	(so Diagram \ref{agat:Eq:ProductDefinitionWithoutUniversalPropertyDiagram} commutes). \\
	
	Show: $h$ bijective. \\
	Construct inverse map $k$:
	\[
	\begin{aligned}
		& k : \text{Mor}_{\mathbf{C}}(T,A) \times \text{Mor}_{\mathbf{C}}(T,B) \to \text{Mor}_{\mathbf{C}}(T,A\times B) \\
		& k((f,g))  := u
	\end{aligned}
	\]
	where $(f,g) \in \text{Mor}_{\mathbf{C}}(T,A) \times \text{Mor}_{\mathbf{C}}(T,B)$
	
	By universal (mapping) property, $\exists \, ! \, u :T\to A\times B$ s.t. $\begin{aligned} & \quad \\ 
	& p_A \circ u=f \\
	& p_B \circ u = g \end{aligned}$.
	
	Let $(f,g) \in \text{Mor}_{\mathbf{C}}(T,A) \times \text{Mor}_{\mathbf{C}}(T,B)$. \\
	Then $(h\circ k)((f\circ g)) = h(u) = (p_A\circ u, p_B \circ u) = (f,g)$. Hence $h\circ k =1$ or i.e. $h\circ k 1_{\text{Mor}_{\mathbf{C}}(T,A) \times \text{Mor}_{\mathbf{C}}(T,B)}$. \\
	
	Let $u\in \text{Mor}_{\mathbf{C}}(T,A\times B)$. \\
	Then $(k\circ h)(u) = k((f, g)) = k((p_A \circ u, p_B \circ u))= u'$ where $u':T\to A\times B$ and \\
	\[
	\begin{aligned} 
	& p_A \circ u' = p_A \circ k((p_A \circ u, p_B \circ u)) = p_A \circ k ((f,g)) = p_A \circ u \quad \, (\text{since $k((f,g)) := u$}) \\
	& p_B \circ u' = p_B \circ k((p_A \circ u, p_B \circ u)) = p_B \circ k ((f,g)) = p_B \circ u \quad \, (\text{since $k((f,g)) := u$}) 
	\end{aligned} 
	\]
	2nd. equality is by universal (mapping) property.

$\Longrightarrow u= u'$ (by component-wise equality), and \\
$k\circ h = 1_{\text{Mor}_{\mathbf{C}}(T,A\times B)}$. \\
$\Longrightarrow h$ isomorphic. \\

$\Longrightarrow$ : \\

Given morphisms
\[
\begin{aligned}
& f:T \to A \\ 
& g:T \to B
\end{aligned}
\]
then $(f, g) \in \text{Mor}_{\mathbf{C}}(T, A) \times \text{Mor}_{\mathbf{C}}(T,B)$.

$\forall \, $ object $T\in \text{Obj}\mathbf{C}$, $\exists \, $ isomorphism $h$ s.t. $\text{Mor}_{\mathbf{C}}(T, A\times B) \cong \text{Mor}_{\mathbf{C}}(T,A) \times \text{Mor}_{\mathbf{C}}(T,B)$, \\
Then $u:=h^{-1}((f, g))$ is a unique morphism, $u\in \text{Mor}_{\mathbf{C}}(T, A\times B)$. \\
Also, by Diagram \ref{agat:Eq:ProductDefinitionWithoutUniversalPropertyDiagram}, 
\[
\begin{aligned}
& \text{Mor}_{\mathbf{C}}(T,p_A)(u) = f = p_{\text{Mor}_{\mathbf{C}}(T,A)} \circ h(u) \\ 
& \text{Mor}_{\mathbf{C}}(T,p_B)(u) = f = p_{\text{Mor}_{\mathbf{C}}(T,B)} \circ h(u) 
\end{aligned}
\]

Now
\[
\begin{aligned}
& \text{Mor}_{\mathbf{C}}(T, p_A)(u) = p_A \circ u \\
& \text{Mor}_{\mathbf{C}}(T, p_B)(u) = p_B \circ u \\
\end{aligned}
\]
(by definition).

So
\[
\begin{aligned} 
& f = p_A \circ u \\ 
& g = p_B \circ u
\end{aligned} 
\]
So universal mapping property is satisfied.
\end{proof}

Thus,

\begin{definition}[Product, generalized, assuming universal property]
	Give $A_i \in \text{Obj}\mathbf{C}$, $i\in I$, category $\mathbf{C}$, \\
	\textbf{product} is object $\prod_{i \in I} A_i \in \text{Obj}\mathbf{C}$, morphisms $p_j : \prod_{i\in I} A_i \to A_j$, $j\in I$, \\
	s.t. $\forall \, $ object $A\in \in \text{Obj}\mathbf{C}$, \\
	$\exists \, !$ morphism $\langle f_i \rangle : A \to \prod_{i\in I} A_i$, \\
	s.t. 
\begin{equation}\label{agat:Eq:GeneralizedProductDiagram}
	\begin{tikzpicture}
\matrix (m) [matrix of math nodes, row sep=5em, column sep=4em, minimum width=3em]
{
	& A  \\ 
	A_j  & \prod_{i\in I} A_i   \\
};
\path[->]
(m-1-2) edge node [left] {$f_j$} (m-2-1)
(m-2-2) edge node [below] {$p_j$} (m-2-1)
;
\path[dashed,->]
(m-1-2)        edge node [right] {$\langle f_i \rangle$} (m-2-2)
;
\end{tikzpicture}  
\end{equation} 	
i.e. $\pi_j \circ \langle f_i \rangle = f_j \, \forall \, i,j \in I$
\end{definition}

\subsubsection{Example: Set always has products}

$\forall \, $ sets $X,Y \in \text{Obj}(\text{Set})$, $\exists \, $ product $X\times Y \in \text{Obj}(\text{Set})$.  

Let $A \in \text{Obj}(\text{Set})$, $\begin{aligned} & \quad \\
& f_1 : A \to X \\
& f_2 : A \to Y \end{aligned}$ \qquad \, Define $\begin{aligned} & \langle f_1, f_2 \rangle : A \to X \times Y  \\
& \langle f_1 ,f_2 \rangle (a) = (f_1(a), f_2(a)) \end{aligned}$

Then $\begin{aligned} & \quad \\
& \pi_1 \circ \langle f_1 , f_2 \rangle (a) = f_1(a) \\
& \pi_2 \circ \langle f_1, f_2 \rangle (a) = f_2(a) \end{aligned}$ \qquad \, $\Longrightarrow \begin{aligned}  & \quad \\
& \pi_1 \circ \langle f_1 , f_2 \rangle = f_1 \\
& \pi_2 \circ \langle f_1,f_2 \rangle = f_2 \end{aligned}$

Suppose $f': A \to X\times Y$ s.t. $\begin{aligned} & \quad \\
& \pi_1 \circ f' = f_1 \\
& \pi_2 \circ f' = f_2 \end{aligned}$

Write $f'(a) = (x,y)$
\[
\begin{aligned}
& f_1(a) = \pi_1 \circ f'(a) = \pi_1(x,y) =x \\ 
& f_2(a) = \pi_2 \circ f'(a) = \pi_2(x,y) =y \\ 
\end{aligned} \qquad \, \Longrightarrow f'(a) = (f_1(a), f_2(a)) = \langle f_1, f_2 \rangle (a)
\]
$\langle f_1, f_2 \rangle $ unique.  

\subsubsection{Example: "Record Types", classes (in computer science, C++, Python)}\label{agat:SubSubSec:RecordTypesProduct}

cf. pp. 174, 5.3.14 "Record Types" Barr and Wells (1998) \cite{agat:BW1998}.

Look at Diagram \ref{agat:Eq:GeneralizedProductDiagram} again.

To allow operations depending on several variables in a functional programming language $L$, \\
e.g. assume $\forall \,$ types $A_j$, $\forall \, j \in I$, language $L$ has record type $\prod_{i\in I} A_i$, and field selectors $p_j$
\[
p_j : \prod_{i\in I} A_i \to A_j
\]
Insist that any data in $\prod_{i\in I} A_i$ be determined completely by those fields $A_j$, $j\in I$, \\
then $\forall \, $ operations $f_j$, $f_j : A \to A_j$, \, $\forall \, j \in I$, \\
there ought to be an unique operation $\lbrace f_i \rangle : A \to \coprod_{i\in I} A_i$ s.t. 
\[
p_j\langle f_i \rangle = f_j, \quad \, \forall \, i, j \in I
\]
This would make $\prod_{i\in I} A_i$ the product of $A_i$'s, with selectors as product projections.

Thus, to say that one can always construct record types in a functional programming language $L$ to say that the corresponding category $C(L)$ has finite products. 


\begin{proposition}
	If product $(A_1 \times \dots \times A_{\mathcal{N}} \xrightarrow{ \pi_i } A_i )_{i\in I}$, if $\exists \,  i_0 \in I$ s.t. $\text{Hom}(A_{i_0}, A_i) \neq \emptyset$, \, $\forall \, i \in I$, \\
	then $\pi_{i_0}$ \emph{retraction }
\end{proposition}

\begin{proof}
	$\forall \, i \in I$, choose $f_i \in \text{Hom}(A_{i_0}, A_i)$ with $f_{i_0} = 1_{A_{i_0}}$.  
	
	Then $\langle f_i \rangle : A_{i_0} \to A_1 \times \dots \times A_{\mathcal{N}}$ is a morphism s.t. 
	\[
	\pi_{i_0} \circ \langle f_i \rangle = f_{i_0} = 1_{A_{i_0}}
	\]
\end{proof}


Ad\'{a}mek, Herrlich, and Strecker (2004) \cite{agat:AHS2004} and their notation) calls a \textbf{sink} what Leinster (2014) \cite{agat:Lein2014} calls a \textbf{cocone}.  
\begin{definition}
	\textbf{sink} $((f_i)_{i\in I}, A) \equiv (f_i , A)_I \equiv (A_i \xrightarrow{ f_i } A)_I$, object $A$, family of morphisms $f_i : A_i \to A$
\end{definition}

For the \emph{coproduct}, consider this enlightening comparision:
\[
\begin{aligned}
\begin{gathered}
\text{ product } ( \prod_{i\in I } A_i, \pi_j )_{j \in I} \\ 
\text{ projection } \pi_j : \prod_{i\in I} A_i \to A_j  \\
\begin{tikzpicture}
\matrix (m) [matrix of math nodes, row sep=5em, column sep=4em, minimum width=3em]
{
	& A  \\ 
	A_j  & \prod_{i\in I} A_i   \\
};
\path[->]
(m-1-2) edge node [above] {$f_j$} (m-2-1)
edge node [right] {$\langle f_i \rangle$} (m-2-2)
(m-2-2) edge node [below] {$\pi_j$} (m-2-1)
;
\end{tikzpicture}  \\
C\xrightarrow{ \langle f, g \rangle } A\times B \\
\prod_{i \in I } f_i  \, , \text{ or if } i = \lbrace 1,2 \rbrace, \, f\times g
\end{gathered} & \qquad \, 
\begin{gathered}
\text{ coproduct } ( \mu_j, \coprod_{i\in I } A_i  )_{j \in I} \\ 
\text{ injection } \mu_j : A_j \to \coprod_{i\in I} A_i  \\
\begin{tikzpicture}
\matrix (m) [matrix of math nodes, row sep=5em, column sep=4em, minimum width=3em]
{
	& A  \\ 
	A_j  & \coprod_{i\in I} A_i   \\
};
\path[->]
(m-2-1) edge node [above] {$f_j$} (m-1-2)
edge node [above] {$\mu_j$} (m-2-2)
(m-2-2) edge node [right] {$[f_i]$} (m-1-2)
;
\end{tikzpicture}  \\
C\xleftarrow{ [ f, g ] } A +  B \\
\coprod_{i \in I } f_i  \, , \text{ or if } i = \lbrace 1,2 \rbrace, \, f +  g
\end{gathered}
\end{aligned}
\]

\subsubsection{Examples (of coproducts)} 

\begin{itemize}
	\item if $(A_i)_I$ pairwise-disjoint family of sets, then $(\mu_j, \bigcup_{i\in I} A_i)_{j\in I}$ is coproduct in $\text{Set}$.  \\
	If $(A_i)_I$ arbitrary set-indexed family of sets, then it can be ``made disjoint'' by pairing each $A_i$ with index $i$, i.e. by working with $A_i \times \lbrace i \rbrace$ rather than $A_i$.  
	
	So $\bigcup_{i\in I } (A_i \times \lbrace i \rbrace)$ disjoint.  Consider 
	\[
	\begin{aligned}
	&  \mu_j : A_j \to \bigcup_{i \in I} A_i \times \lbrace i \rbrace \\ 
	&  \mu_i(a) = (a,j)
	\end{aligned}
	\]
	$(\mu_j, \bigcup_{i\in I} A_i \times \lbrace i \rbrace )_{j\in I}$ is a coproduct in $\text{Set}$.  
	
	Indeed, given $\begin{aligned} & \quad \\
	& f_j:A_j \to A \\
	& f_j(a) \in A \end{aligned}$, 
	\[
	\begin{aligned}
	& [f_i] : \coprod_{i\in I} A_i \times \lbrace i \rbrace \to A \\ 
	& [f_i]\circ \mu_j = f_j 
	\end{aligned}
	\]
	where
	\[
	f_j(a) = [f_i]\circ \mu_j(a) = [f_i](a,j) = f_j(a)
	\]
	\item  $\text{Top}$ coproducts are ``topological sums''; they're ``concrete'' coproducts (Ad\'{a}mek, Herrlich, and Strecker (2004) \cite{agat:AHS2004})
	\item $\text{Vec}$ (nonconcrete) coproducts called \emph{direct sums}
	
	direct sum $\bigoplus_{i\in I} A_i$ of vector spaces $A_i$ is subspace of direct product $\prod_{i\in I} A_i$ \\
	consisting of all elements $(a_i)_{i\in I}$ with finite carrier (i.e. $\lbrace i \in I | a_i \neq 0 \rbrace$ is finite), \\
	\qquad injections \[
	\begin{aligned} & \quad \\
	& \mu_j : A_j \to \bigoplus_{i\in I } A_i  \\
	& \mu_j(a) = (a_i)_{i\in I} \text { with } a_i = \begin{cases} a & \text{ if } i = j  \\ 0 & \text{ if } i \neq j \end{cases}
	\end{aligned}
	\]
	
	\item $\text{Grp}$ has nonconcrete coproducts, ``free products''
\end{itemize}

\section{Naturality; Natural transformations}

\begin{definition}[Natural Transformation]\label{agat:Def:NaturalTransformation}
	If functors $F,G$, categories $\mathbf{C}, \mathbf{D}$, 
	\[
	\begin{aligned}
	& F:\mathbf{C} \to \mathbf{D} \\ 
	& G:\mathbf{C} \to \mathbf{D} 
	\end{aligned}
	\]
	the \textbf{natural transformation} $\eta$ from $F$ to $G$ is a family of morphisms s.t. \\
	
	\begin{enumerate} 
		\item $\forall \, X \in \text{Obj}\mathbf{C}$, $\eta_X : F(X) \to G(X)$ s.t. $F(X), G(X) \in \text{Obj}\mathbf{D}$. morphisms $\eta_X \in \text{Hom}_{\mathbf{D}}(F(X), G(X)) \subset \text{Mor}\mathbf{D}$ is called \textbf{component} of $\eta$ at $X$.
		\item components must be s.t. $\forall \, $ morphism $f: X \to Y$ in $\mathbf{C}$, i.e. $f\in \text{Hom}_{\mathbf{C}}(X, Y) \in \text{Mor}\mathbf{C}$, 
		\begin{equation}
		\eta_Y \circ F(f) = G(f) \circ \eta_X
		\end{equation} 
		i.e. 
		\[
		       \begin{tikzpicture}
		\matrix (m) [matrix of math nodes, row sep=5em, column sep=5em, minimum width=2em]
		{
			       F(X) & F(Y) \\
			       G(X) & G(Y) \\
			};
		\path[->]
		(m-1-1) edge node [above] {$F(f)$} (m-1-2)
		edge node [auto] {$\eta_X$} (m-2-1)
		(m-1-2) edge node [auto] {$\eta_Y$} (m-2-2)
		(m-2-1) edge node [auto] {$G(f)$} (m-2-2);
		\end{tikzpicture}
		\]
		
		If both $F,G$ contravariant, 
		\[
		\begin{tikzpicture}
		\matrix (m) [matrix of math nodes, row sep=5em, column sep=5em, minimum width=2em]
		{
			       F(X) & F(Y) \\
			       G(X) & G(Y) \\
			};
		\path[->]
		(m-1-2) edge node [above] {$F(f)$} (m-1-1)
		(m-1-1) edge node [auto] {$\eta_X$} (m-2-1)
		(m-1-2) edge node [auto] {$\eta_Y$} (m-2-2)
		(m-2-2) edge node [auto] {$G(f)$} (m-2-1);
		\end{tikzpicture}
		\]
		\end{enumerate} 
	\end{definition}
cf. \href{https://en.wikipedia.org/wiki/Natural_transformation}{Wikipedia, "Natural transformation"}

Notation: natural transformation from $F$ to $G$, $\eta:F \to G$, i.e. family of morphisms $\eta_X : F(X) \to G(X)$ is natural in $X$.

cf. Ad\'{a}mek, Herrlich, and Strecker (2004) \cite{agat:AHS2004}.  \\

From pp. 16, Sec. 14 "Natural Transformations", Mac Lane (1978) \cite{agat:MacL1978}, \\

$\forall \,$ arrow $f: C \to C'$ in $C$, i.e. notation: $\forall \, $ morphism $f:X \to Y$ in $\mathbf{C}$, 

Mac Lane's notation:
\[
\begin{gathered}
\begin{tikzcd} 
c \arrow{d}{f} \\
c'
\end{tikzcd} \qquad \, 
\begin{tikzcd}
Sc \arrow{d}{Sf} \arrow{r}{\tau c} & Tc \arrow{d}{Tf} \\
Sc' \arrow{r}{\tau c'} & Tc'
\end{tikzcd}
\end{gathered}
\]

My notation:
\[
\begin{gathered}
\begin{tikzcd} 
X \arrow{d}{f} \\
Y
\end{tikzcd} \qquad \,
\begin{tikzcd}
FX \arrow{d}{Ff} \arrow{r}{\eta_X} & GX \arrow{d}{Gf} \\
FY \arrow{r}{\eta_Y} & GY
\end{tikzcd}
\end{gathered}
\]

We say that $\tau_c : Sc \to Tc$ is natural in $C$; $\eta_X : FX \to GX$ is natural in $X$. \\

If we think of functor $S$ as giving a picture in $B$ of (all the objects, arrows of) $C$, then \\
natural transformation $\tau$ is set of arrow $s$ mapping (or, translating) the picture $S$ to picture $T$, with all squares (and parallelograms!) like that above commutative:

Mac Lane's notation:
\[
\begin{gathered}
\begin{tikzcd}
a \arrow{dr}{f} \arrow{dd}{h} & \\
& b \arrow{dl}{g} \\ 
c & 
\end{tikzcd} \qquad \, 
\begin{tikzcd}
Sa \arrow{dr}{Sf} \arrow{dd}{Sh} \arrow{rr}{\tau a} & & Ta \arrow[dd, crossing over]{} \arrow{dr}{Tf} &  \\
& Sb \arrow{dl}{Sg} \arrow{rr}[swap, pos=0.2]{\tau b} & & Tb \arrow{dl}{Tg} \\ 
Sc \arrow{rr}{\tau c} & & Tc & 
\end{tikzcd} 
\end{gathered}
\]

My notation:
\[
\begin{gathered}
\begin{tikzcd}
X \arrow{dr}{f} \arrow{dd}{g} & \\
& Y \arrow{dl}{h} \\ 
Z & 
\end{tikzcd} \qquad \,
\begin{tikzcd}
FX \arrow{dr}{Ff} \arrow{dd}{Fg} \arrow{rr}{\eta_X} & & GX \arrow[dd, crossing over]{} \arrow{dr}{Gf} &  \\
& FY \arrow{dl}{Fh} \arrow{rr}[swap, pos=0.2]{\eta_Y} & & GY \arrow{dl}{Gh} \\ 
FZ \arrow{rr}{\eta_Z} & & GZ & 
\end{tikzcd} 
\end{gathered}
\]

\emph{natural equivalence} or i.e. natural isomorphism - natural transformation $\tau$ with $\forall \, $ component $\tau c$ invertible in $B$, $\tau : S \cong T$.  \\
notation: $\eta$ with $\forall \, $ component $\eta_X$ invertible in $\mathbf{D}$; $\eta : F \cong G$. \\
$\Longrightarrow $ inverses $(\tau c)^{-1}$ in $B$ are components of natural isomorphism $\tau^{-1} :T \to S$. \\
notation: inverses $(\eta_X)^{-1}$ in $\mathbf{D}$ are components of natural isomorphism $\eta^{-1}: G\to F$. 


\subsection{Examples of Natural transformations}

e.g. Hurewicz homomorphism $\pi_n(X) \to H_n(X)$ $\forall \, $ topological space $X$ is a natural transformation from $n$th homology functor. \\
$\pi_n : \textbf{Top} \to \mathbf{Grp}$ to $n$th homology functor $H_n : \textbf{Top} \to \textbf{Grp}$.

e.g. \textbf{determinant} $\det$ is a natural transformation (cf. Mac Lane (1978) \cite{agat:MacL1978}). \\
Let $\det_KM = $ determinant of $n\times n$ matrix $M$, with entries in commutative ring $K$, while \\
$K^* \equiv $ group of units (invertible elements) of $K$ \\

Thus $M$ non-singular when $\det_KM$ is a unit, and $\det_K$ is a morphism $GL_nK \to K^*$ of groups. \\
Because $\det$ is defined by same formula $\forall \,$ ring $K$, each morphism $f: K \to K'$ of commutative rings leads to 
\[
\begin{tikzcd}
GL_nK \arrow{d}{GL_nf} \arrow{r}{\det_K} & K^* \arrow{d}{f^*} \\
GL_nK' \arrow{r}{\det_{K'}} & (K')^*
\end{tikzcd}
\]

$\Longrightarrow \det : GL_n \to ()^*$ is natural between 2 functors $\textbf{CRng} \to \textbf{Grp}$

TODO: Work out examples of natural transformations involving group commutators on pp. 17 of Mac Lane (1978) \cite{agat:MacL1978}.

e.g. category $\textbf{Finord} = $ all finite ordinal numbers $n$ \\
category $\textbf{Set}_f = $ all finite sets (in some universe $U$) \\ 
$\forall \, $ ordinal $n = \lbrace 0 , 1 , \dots n-1 \rbrace$ is a finite set, so inclusion $S$ is a functor $S: \textbf{FinOrd} \to \textbf{Set}_f$ \\
$\forall \, $ finite set $X$ determines ordinal number $n = \# X$, number of elements in $X$ \\
$\forall \, X$, choose bijection $\theta_X:X \to \# X$ \\
$\forall \, $ function $f:X \to Y$, $X,Y$ finite sets, \\
Choose corresponding function $\# f : \# X \to \# Y$ between ordinals by $\# f = \theta_Y f \theta_X^{-1}$, so that 
\begin{equation}\label{agat:Eq:NaturalIsomorphismOrdinalNumbersFiniteSets}
\begin{tikzcd}
X \arrow{d}{f} \arrow{r}{\theta_X} & \# X \arrow{d}{\# f} \\
Y \arrow{r}{\theta_Y} & \# Y
\end{tikzcd}
\end{equation}

$\Longrightarrow \# $ a functor $\# : \textbf{Set}_f \to \textbf{FinOrd}$ \\
If $X$ itself an ordinal number, take $\theta_X$ to be an identity. \\
$\Longrightarrow \# \circ S$ is identity functor $I'$ of $\textbf{FinOrd}$ \\
$S \circ \#$ is not identity functor $I : \textbf{Set}_f \to \textbf{Set}_f$ because \\
$S \circ \#$ sends each finite set $X$ to a special finite set - ordinal number $n$ with same number of elements as $X$. \\
However $\theta: I \to S \#$ is a natural isomorphism by square diagram above (Eq. \ref{agat:Eq:NaturalIsomorphismOrdinalNumbersFiniteSets}).  
\[
I \cong S\circ \# , \quad \, I' = \# \circ S 
\]

equivalence between categories $\mathbf{C}, \mathbf{D} =$ \\
pair of functors $\begin{aligned} & \quad \\ 
& S : \mathbf{C} \to \mathbf{D} \\ 
& T : \mathbf{C} \to \mathbf{D} \end{aligned}$ and \\
natural isomorphisms $I_C \cong T \circ S$, $I_D \cong S \circ T$. 





\section{Limits}

\subsection{Pullback}

\begin{definition}
	For some category $\mathbf{A}$, and for 
	
	\begin{tikzpicture}
	\matrix (m) [matrix of math nodes, row sep=5em, column sep=5em, minimum width=2em]
	{
		& Y \\
		X & Z \\
	};
	\path[->]
	(m-1-2) edge node [right] {$t$} (m-2-2)
	(m-2-1) edge node [below] {$s$} (m-2-2);
	\end{tikzpicture}
	
	$X,Y,Z \in \text{Obj}\mathbf{A}$.  
	
	$\begin{aligned} & \quad \\
	& s : X \to Z \\
	& t: Y \to Z \end{aligned}$ ; \qquad \, $s,t \in \text{Mor}\mathbf{A}$
	
	Then the \textbf{pullback} or ``pullback square'' consists of $P \in \text{Obj}\mathbf{A}$, $\begin{aligned} & \quad \\
	& \pi_1 : P \to X \\
	& \pi_2: P \to Y \end{aligned}$ s.t. 
	
	\begin{tikzpicture}
	\matrix (m) [matrix of math nodes, row sep=5em, column sep=5em, minimum width=2em]
	{
		P & Y \\
		X & Z \\
	};
	\path[->]
	(m-1-1) edge node [above] {$\pi_2$} (m-1-2)
	edge node [auto] {$\pi_1$} (m-2-1)
	(m-1-2) edge node [auto] {$t$} (m-2-2)
	(m-2-1) edge node [auto] {$s$} (m-2-2);
	\end{tikzpicture}
	
	commutes and s.t. $\forall \, $ commutative square in $\mathbf{A}$ 
	
	\begin{tikzpicture}
	\matrix (m) [matrix of math nodes, row sep=5em, column sep=5em, minimum width=2em]
	{
		A & Y \\
		X & Z \\
	};
	\path[->]
	(m-1-1) edge node [above] {$f_2$} (m-1-2)
	edge node [auto] {$f_1$} (m-2-1)
	(m-1-2) edge node [auto] {$t$} (m-2-2)
	(m-2-1) edge node [auto] {$s$} (m-2-2);
	\end{tikzpicture}
	
	then $\exists \, ! \, f: A \to P$ s.t. 
	
	\begin{tikzpicture}
	\matrix (m) [matrix of math nodes, row sep=5em, column sep=5em, minimum width=2em]
	{
		A &   & \\ 
		& P & Y \\
		& X & Z \\
	};
	\path[->]
	(m-2-2) edge node [above] {$\pi_2$} (m-2-3)
	edge node [auto] {$\pi_1$} (m-3-2)
	(m-2-3) edge node [auto] {$t$} (m-3-3)
	(m-3-2) edge node [auto] {$s$} (m-3-3)
	(m-1-1) edge node [above] {$f_2$} (m-2-3)
	edge node [above] {$f$} (m-2-2)
	edge node [below] {$f_1$} (m-3-2);
	\end{tikzpicture}
	
	
\end{definition}

Simmons (2011) \cite{agat:Simm2011}

cf. 3.3 \emph{Some less simple functors}, pp. 79, of Simmons (2011) \cite{agat:Simm2011}

\section{Monads}

cf. Ch. 6, Monads and Algebras, starting on pp. 137 of Mac Lane (1971) \cite{agat:MacL1978}.

\subsection{Monads in a Category}

\subsubsection{Development of a Monad}

Let's follow the development of a Monad according to Mac Lane (1971) \cite{agat:MacL1978} starting on pp. 137 in Sec. 1. "Monads in a Category."

$\forall \, $ endofunctor $T: X \to X$, $T$ has composites $T^2 = T\circ T : X \to X$, $T^3 = T^2 \circ T:X \to X$ in Mac Lane's notation, $X$ is a category. In our notation, $\mathbf{C}$ denotes the category. And so for endofunctor $T: \mathbf{C} \to \mathbf{C}$, $T^2 = T\circ T: \mathbf{C} \to \mathbf{C}$, $T^3 = T^2 \circ T: \mathbf{C} \to \mathbf{C}$. 

Recall from Def. \ref{agat:Def:Endofunctor}, that for endofunctor $T:\mathbf{C} \to \mathbf{C}$, the \emph{object map} $T: \text{Obj}{\mathbf{C}} \to \text{Obj}{\mathbf{C}}$ acts just like a morphism in $\text{Mor}{\mathbf{C}}$! So $T: X \to T(X)$ for $X, T(X) \in \text{Obj}{\mathbf{C}}$. Also, recall that the \emph{morphism map} $T: \text{Mor}_{\mathbf{C}} \to \text{Mor}_{\mathbf{C}}$ acts like the usual morphism map, such that $Tf : T(X) \to T(Y)$ if $f:X\to Y$, for $X,Y \in \text{Obj}{\mathbf{C}}$. \\

If $\mu:T^2 \to T$ is a \emph{natural transformation}, then recall the definition of the natural transformation, Def. \ref{agat:Def:NaturalTransformation}, and apply that onto $\mu$. \\

$\forall \, X \in \text{Obj}{\mathbf{C}}$, $\mu_X: T^2(X) \to T(X)$ with $T^2(X), T(X) \in \text{Obj}{\mathbf{C}}$. $\mu_X$ is a morphism $\in \text{Hom}_{\mathbf{C}}(T^2(X), T(X))$ (!!), as well as being a component of $\mu$ at $X$. \\

$\forall \,$ morphism $f: X\to Z$ in $\mathbf{C}$, $X, Z \in \text{Obj}{\mathbf{C}}$ (the notation of using $Z$ to emphasize that $f$ can map $X$ to \emph{any} arbitrary object in $\text{Obj}{\mathbf{C}}$), 
\[
\mu_Z \circ T^2f = Tf \circ \mu_X 
\]
i.e. 
\[
\begin{tikzcd}
	T^2(X) \arrow{d}{T^2f} \arrow{r}{\mu_X} & T(X) \arrow{d}{Tf} \\
	T^2(Z) \arrow{r}{\mu_Z} & T(Z)
\end{tikzcd}
\]

Let $T\mu : T^3 \to T^2$ \emph{denote} a natural transformation with components $(T\mu)_X = T(\mu_X) : T^3(X) \to T^2(X)$, \\
while $\mu T: T^3\to T^2$ denotes another natural transformation with components $(\mu T)_X  = \mu_{T(X)}$.  

$T\mu$, $\mu T$ are "horizontal" composites in the sense of Sec. II. 5, "The Category of All Categories". (TODO: Understand what Mac Lane (1971) \cite{agat:MacL1978} means when he says this on pp. 137). \\

Let's jump the gun and consider what natural transformation $\eta : 1_{\mathbf{C}} \to T$ means with respect to the definition of a natural transformation, Def. \ref{agat:Def:NaturalTransformation}. $1_{\mathbf{C}}$ is the identity functor, mapping everything to itself. \\

$\forall \, X \in \text{Obj}{\mathbf{C}}$, $\eta_X: 1_{\mathbf{C}}(X) \to T(X)$ with $1_{\mathbf{C}}(X) = X, T(X) \in \text{Obj}{\mathbf{C}}$. $\eta_X$ is a morphism $\in \text{Hom}_{\mathbf{C}}(1_{\mathbf{C}}(X), T(X)) = \text{Hom}_{\mathbf{C}}(X, T(X))$ (!!!), as well as being a component of $\eta$ at $X$. Not only that, it applies $T$ onto $X$! \\

$\forall \,$ morphism $f: X\to Z$ in $\mathbf{C}$, $X, Z \in \text{Obj}{\mathbf{C}}$, 
\[
\eta_Z \circ 1_{\mathbf{C}}f = \eta_Z \circ f = Tf \circ \eta_X 
\]
i.e. 
\[
\begin{tikzcd}
1_{\mathbf{C}}(X) \arrow{d}{1_{\mathbf{C}}f} \arrow{r}{\eta_X} & T(X) \arrow{d}{Tf} \\
1_{\mathbf{C}}(Z) \arrow{r}{\eta_Z} & T(Z)
\end{tikzcd} \Longrightarrow 
\begin{tikzcd}
X \arrow{d}{f} \arrow{r}{\eta_X} & T(X) \arrow{d}{Tf} \\
Z \arrow{r}{\eta_Z} & T(Z)
\end{tikzcd}
\]

\subsubsection{Definitions and Properties of a Monad}

\begin{definition}[Monad]\label{agat:Def:Monad}
\textbf{monad} $T = \langle T, \eta, \mu \rangle$ in category $\mathbf{C}$ consists of functor $T: \mathbf{C} \to \mathbf{C}$ and 2 natural transformations
\begin{equation}
\begin{aligned}
& \eta : 1_{\mathbf{C}} \to T \\ 
& \mu : T^2 \to T
\end{aligned}
\end{equation}	
s.t. 
\begin{equation}\label{agat:Eq:MonadAssociativeDiagram}
\begin{tikzcd}
T^3 \arrow{d}{\mu T} \arrow{r}{T\mu} & T^2 \arrow{d}{\mu} \\
T^2 \arrow{r}{\mu} & T
\end{tikzcd}
\end{equation}
being like an "associative" diagram, and 
\begin{equation}\label{agat:Eq:MonadLeftRightUnitLawsDiagram}
\begin{tikzcd}
	1T \arrow[equal]{d} \arrow{r}{\eta T} & T^2 \arrow{d}{\mu} & T1 \arrow{l}[swap]{T \eta} \arrow[equal]{d} \\
	T \arrow[equal]{r} & T & T \arrow[equal]{l}
\end{tikzcd}
\end{equation}
being a diagram that could be said to express the left and right unit laws.
\end{definition} 

Formally, the monad definition is like that of a monoid $M$ in $\mathbf{Sets}$: \\
set $M$ of elements of monad \quad $\to $ \, replaced by endofunctor $T : X \to X$ \\
Cartesian product $\times$ of 2 sets \quad $\to $ \, composite of 2 functors \\
binary operation $\mu : M \times M \to M$ of multiplication \quad $\to $ \, transformation $\mu : T^2 \to T$ and \\
unit (identity) element $\eta:1 \to M$ \quad $\to $ \, $\eta : 1_{\mathbf{C}} \to T$. \\

Thus, call \\
$\eta$ the \textbf{unit} of monad $T$. The first commutative diagram, Eq. \ref{agat:Eq:MonadAssociativeDiagram}, is the associative law for a monad, and \\
$\mu$ the \textbf{multiplication} of monad $T$. The 2nd. and 3rd. diagrams, Eq. \ref{agat:Eq:MonadLeftRightUnitLawsDiagram} expresses the left and right \emph{unit laws}.

\begin{proposition}[Coherence conditions]\label{agat:Prop:MonadCoherenceConditions}
	\begin{equation}
	\begin{aligned}
	& \mu \circ T \mu = \mu \circ \mu T \quad (\text{as natural transformations $T^3 \to T$, $\mu \circ T\mu$ and $\mu \circ \mu T$, that is}) \\ 
	& \mu \circ T \eta = \mu \circ \eta T = 1_T \quad (\text{as natural transformations $T \to T$, $1_T \equiv$ identity functor}) 
	\end{aligned}
	\end{equation}
	
	For the last statement, it is equal to commutative diagram Eq. \ref{agat:Eq:MonadLeftRightUnitLawsDiagram}, since it can be redrawn as follows:
\[
\begin{tikzcd}
	T \arrow{d}{ T \eta} \arrow{r}{\eta T} \arrow[equal]{dr} & T^2 \arrow{d}{\mu} \\
	T^2 \arrow{r}{\mu} & T
\end{tikzcd}	
\]
	
\end{proposition}

\subsection{Monads and Adjoints}

cf. Ch. 6, Monads and Algebras, pp. 138 of Mac Lane (1971) \cite{agat:MacL1978}.
\begin{proposition}
	adjunction $\langle F, G, \eta, \epsilon \rangle : X \to A$, $\exists \, $ monad in category $X$ (notation $\langle F,G, \eta, \epsilon \rangle : \mathbf{C} \to \mathbf{D}$, $\exists \, $ monad in category $\mathbf{C}$). 
	
	Specifically, 2 functors $F:\mathbf{C} \to \mathbf{D}$, $G:\mathbf{D} \to \mathbf{C}$ have composite $T = GF$, an endofunctor, \\
	unit $\eta$ of adjunction is natural transformation $\eta : 1 \to T$, \\
	counit $\epsilon :FG \to 1_A$ of adjunction yields by horizontal composition; a natural transformation
	\[
	\mu = G\epsilon F: GFGF \to GF = T
	\] 
\end{proposition}

TODO: adjunction.

\subsection{Applications of Monads: Correspondence between Monads in Computer Science and Monads in Category Theory}

cf. \href{https://ncatlab.org/nlab/show/monad+%28in+computer+science%29}{monad (in computer science), nLab}

For monads in computer science, \\
maps type $X$ to new type $T(X)$, \\
equipped with rule for composing 2 functions of form $f:X \to T(Y)$ (called Kleisli functions), $g:Y \to T(Z)$ to function 
\[
g\circ f : X \to T(Z) \quad \text{(their Kleisi composition)}
\]
associative in evident sense, unital with respect to given unit function called $\text{pure}_X: X \to T(X)$ to be thought of as taking a value to the pure computation that simply returns that value. \\

e.g. when monad $T(-)$ forms product types $T(X) = X \times Q$, some fixed type $Q$ that carries structure of monoid. \\
Then Kleisi function $f:X \to Y \times Q$ is a function $X\to Y$  that produces as a side effect output of type $Q$. \\
Kleisi composition of $\begin{aligned} &  \quad \\
	& f:X\to Y\times Q \\
	& g: Y \to Z\times Q\end{aligned}$ \, not only evaluates the 2 programs in sequence, but also combines their $Q$-output using monoid operation of $Q$, e.g.
	
	if $\begin{aligned} & \quad \\ 
	& fx = (y, q) \\ 
	& gy = (z,q')\end{aligned}$  \, then $(g\circ f)(x) = (z, qq')$ \\
	
Let syntactic category be denoted $\mathbf{C}$. \\

$\text{Obj}{\mathbf{C}} \ni X$ \quad \, $\to $ \, types of the (programming) language \\
$\text{Hom}_{\mathbf{C}}(X,Y)$ \quad \, $\to $ \, "programs" that takes value of type $X$ as input and returns value of type $Y$ \\

endofunctor $T: \mathbf{C} \to \mathbf{C}$ \\
object map $T:X \to T(X)$ sends each type $X\in \text{Obj}{\mathbf{C}}$ to another type $T(X) \in \text{Obj}{\mathbf{C}}$.  \\

unit natural transformation $\epsilon:1_{\mathbf{C}} \to T$ provides $\forall \,$ type $X$ a component morphism $\text{pure}_X :X \to T(X)$ (notation, \href{https://ncatlab.org/nlab/show/monad+%28in+computer+science%29}{ncat} uses the notation $\epsilon$; we will use $\eta$ for the "unit" natural transformation of a monad):
	\[
		\begin{aligned}
		&  \eta : 1_{\mathbf{C}} \to T \\ 
		& \eta_X : X\to T(X) \equiv \text{pure}_X : X \to T(X)
		\end{aligned}	
	\]

multiplication natural transformation $\mu :T^2 \to T$ provides $\mu_X : T^2(X) \to T(X)$ which induces Klesli composition by formula
\begin{equation}
\begin{aligned} 
	(g\circ f) & := (Y \xrightarrow{g} T(Z)) \circ_{\text{Klesli}} (X\xrightarrow{f} T(Y)) \\
	& := X \xrightarrow{f} T(Y) \xrightarrow{T(g)} T(T(Z)) \xrightarrow{ \mu(Z) } TZ
\end{aligned} 
\end{equation}
$T(g)$ uses fact that $T(-)$ is functor morphism $\mu(Z)$ is 1 that implements "$T$-computation".

Now \href{https://ncatlab.org/nlab/show/monad+%28in+computer+science%29}{nLab} says this about the "bind" operation: \\
	"bind" \quad $\leftrightarrow$ \, multiplication on monad $M$, \\
	\begin{equation}\label{agat:Eq:BindAsStrongMonad}
	\begin{gathered} 
	MA \to (MB)^A \to MB \text{ equivalent to map } MA \times MB^A \to MB \\ 
	 MA \times MB^A \xrightarrow{ \text{ strength } } M(A\times MB^A) \xrightarrow{ M_{\text{eval}_{A, MB} } } MMB \xrightarrow{ mB } MB
	\end{gathered} 
	\end{equation}
To try to make sense of Eq. \ref{agat:Eq:BindAsStrongMonad}, consider the following: 
\[
\begin{gathered}
\text{Recall for natural transformation $\eta$, the component $\eta_X$: } \mu_X : T^2(X) \to T(X) \text{ and for some morphism $f \in \text{Mor}{\mathbf{C}}$}, \\
	f:X\to Z \\
	Tf:T(X) \to T(Z)
\end{gathered} 
\]
So if $f:X\to T(Y)$ instead, 
\[
\begin{gathered}
Tf:T(X) \to T^2(Y) \\ 
(\mu_Y \circ Tf) : T(X) \to T(Y)
\end{gathered} 
\]
Form a Cartesian product of $T(X)$ and $\text{Hom}_{\mathbf{C}}(X,T(Y)) \ni f$ as inputs: $T(X) \times \text{Hom}_{\mathbf{C}}(X, T(Y))$. \\

Then I suggest that Eq. \ref{agat:Eq:BindAsStrongMonad} is translated as follows:
\begin{equation}
T(X) \times \text{Hom}_{\mathbf{C}}(X, T(Y)) \xrightarrow{ \text{ apply } T } T( T(X) \times \text{Hom}_{\mathbf{C}}(X, T(Y))) \xrightarrow{ Tf} T^2(Y) \xrightarrow{ \mu_Y } T(Y) 
\end{equation}
where apply $T$ is on $f\in \text{Hom}_{\mathbf{C}}(X,T(Y))$ in order to form $Tf$.

\subsubsection{More on Klesli composition}

cf. \href{https://ncatlab.org/nlab/show/Kleisli+category}{Kleisli category, nLab}.

\begin{proposition}\label{agat:Prop:KlesliCompositionForm}
	Klesli composition of $g\circ f$ of $\begin{aligned} & \quad \\ 
	& f: X\to T(Y) \\ 
	& g: Y \to T(Z) \end{aligned}$ \, is given by 
	\[
	X \xrightarrow{f} TY \xrightarrow{ Tg} T^2Z \xrightarrow{\mu_Z} TZ 
	\]
	i.e. 
	\begin{equation}
		\boxed{ g \circ_{\text{Klesli}} f = \mu_Z \circ Tg \circ f }
	\end{equation}
	with $ \mu_Z \circ Tg \circ f : X \to T(Z)$
\end{proposition}

\begin{proof}
	Fullness of functor $T$: \\
	
	$\forall \, $ morphism $g: T(X) \to T(Y)$ has antecedent the composition
	\[
	X\xrightarrow{ \eta_X} T(X) \xrightarrow{g} T(Y) 
	\]
	i.e. $g\circ \eta_X : X \to T(Y)$. \\
	
Consider $T(g \circ \eta_X) = Tg \circ T \eta_X : T(X) \to T^2(Y)$ for 
\[
\begin{aligned} 
	& T\eta_X : T(X) \to T^2(X) \\ 
	& Tg : T^2(X) \to T^2(Y) \\ 
	& \mu_Y : T^2(Y) \to T(Y) 
\end{aligned} 
\]
So 
\[
\mu_Y \circ Tg \circ T\eta_X = \mu_Y \circ T(g\circ \eta_X) : T(X) \to T(Y)
\]
Because $g$ is a morphism of algebras (TODO: understand what this means, from \href{https://ncatlab.org/nlab/show/Kleisli+category}{Kleisli category, nLab}), 
\[
\mu_Y \circ Tg \circ T\eta_X = g\circ \mu_X \circ T\eta_X = g
\]
where $\mu_Y \circ Tg = g\circ \mu_X$ uses natural transformation definition \ref{agat:Def:NaturalTransformation} for the property $\mu_Z \circ T^2 f = Tf \circ \mu_X$ (TODO: Show $g=Tf$ is of this form), \\
and $\mu \circ T\eta = \mu \circ \eta T = 1_T$ (coherence condition from Prop. \ref{agat:Prop:MonadCoherenceConditions}).  \\

Faithfulness: \\

If $\mu_Y \circ Tf = \mu_Y \circ Tg$, then precompose by $\eta_X$ yields 
\[
\mu_Y \circ Tf \circ \eta_X = \mu_Y \circ  \eta_{f(X)} \circ f = \mu_Y \circ \eta_{T(Y)} \circ f = f
\]
where \\
$Tf\circ \eta_X = \eta_{T(Y)} \circ f$ from natural transformation def. \ref{agat:Def:NaturalTransformation} was used, and $\mu \circ \eta T = 1_T$ (coherence condition) was also used, for $f:X \to T(Y)$. \\

Similarly for $g$. $\mu_Y \circ Tf = \mu_Y \circ Tg \Longrightarrow f=g$.
\end{proof} 

Notice that the multiplication natural transformation $\mu:T^2 \to T$ was needed to induce the Klesli composition.

\subsubsection{Dictionary between Monads in Computer Science and Monads in Category Theory}

\[
\boxed{ \texttt{join} \quad \text{ is } \, \mu_X : T^2(X) \to T(X)}
\]

\[
\boxed{ \texttt{return} \quad \text{ is } \, \eta_X: X \to T(X) }
\]

\[
\boxed{ \texttt{fmap} \quad \text{ is the morphism map (arrow map) for functor } F}
\]

\[
\boxed{ \texttt{bind} \quad \text{ is } \, \mu_Y \circ Tf : T(X) \to T(Y) \text{ but ... }}
\]
\texttt{bind} is more nuanced in that, from the Haskell and functional programming perspective, it takes in \emph{two} inputs. \\

One of them is an element $x$ in the object $X$, $x \in X$. The other is a "transformation of type $X \to T(Y)$." I believe that would mean a morphism $f: X \to T(Y)$.  \\

What \v{C}uki\'{c} in \v{C}uki\'{c} (2018) \cite{agat:Cuki2018} calls \texttt{transform} is the \emph{morphism map} (arrow map) for a functor $F$. However, note that when it comes to concrete implementation, the concrete implementation applies to \emph{instances} of a type $X \in \text{Obj}{\mathbf{C}}$. For a given type $X$, e.g. \verb|int| (the integer type in C or C++), there are many values that will have that type. It would be good that we work out how the definitions, theorems, and commutative diagrams apply in this case.



I don't think there's a great way to reconcile what Haskell and functional programming programmers mean by \texttt{bind} with the corresponding concept in category theory. 

\subsection{Examples of Monads in Computer Science}


\subsubsection{Reader Monad; Function Monad; Environment Monad} $[W, -] : \mathbf{C}\to \mathbf{C}$.

Given $W \in \text{Obj}{\mathbf{C}}$, ($W$ is a type), 
\[
[W, - ] = (W \to (-))
\]
where $[-, -] : \mathbf{C}^{\text{op}} \times \mathbf{C} \to \mathbf{C}$.

$W \mapsto $ internal hom out of $W$.

\[
\begin{aligned} 
& \eta_X : X \to [W, -](X) \\ 
& \eta_W: W \to [W, -](W) \\ 
& \eta_W : w \mapsto [W,-](W)(w) \text{ \, value to constant functions with that value } \\ 
& \Delta : X \in \text{Obj}{\mathbf{C}} \xrightarrow{ (1,1) } X \times X \\ 
& \Delta : x \mapsto (w,w)
\end{aligned} 
\]

"read in " state of type $W$ (an environment)

\href{https://en.wikipedia.org/wiki/Monad_(functional_programming)#State_monads}{wikipedia} calls this the Environment Nomad. \\

\texttt{reader}: $T \to E \to T = t\mapsto e \mapsto t$ corresponds to $\eta_X : X \to T(X)$ where $T(X) = \text{Hom}_{\mathbf{C}}(E,X)$ s.t. $\eta_X$ sends values $x\in X $ to constant functions (on $E$) to value $x$. \\

Therefore, rewrite the unit component at $X$, i.e. \texttt{reader}, as \\
\begin{equation}
\begin{gathered}
\begin{aligned}
& X \to E \to X \\
& x \mapsto e \mapsto x
\end{aligned} \qquad \, 
\begin{aligned}
	& \eta_X : X \to \text{Hom}_{\mathbf{C}}(E,X) \\ 
	& \eta_X : x \mapsto (e\mapsto x)
\end{aligned}
\end{gathered}
\end{equation}
Thus, the \emph{unit} component at $X$, called \texttt{unit} or \texttt{return} by programmers, truly maps a value $x \in X$ to a \emph{constant function}:
\begin{equation}
\eta_X : x \mapsto (e \mapsto x) \in \text{Hom}_{\mathbf{C}}(E, X)
\end{equation}


\texttt{bind}: $(E\to T) \to (T\to E \to T') \to E \to T' = r\mapsto f \mapsto e \mapsto f(re)e$ corresponds to $\mu_Y \circ Tf:T(X) \to T(Y)$ where $T(X) = \text{Hom}_{\mathbf{C}}(E,X)$, $T(Y) = \text{Hom}_{\mathbf{C}}(E,Y)$, and \\
$f:X\to T(Y)$ or $f:X\to \text{Hom}_{\mathbf{C}}(E,Y)$ \\
$Tf : T(X) \to T^2(Y)$ or $Tf : \text{Hom}_{\mathbf{C}}(E,X) \to \text{Hom}_{\mathbf{C}}(E,T(Y)) = \text{Hom}_{\mathbf{C}}(E, \text{Hom}_{\mathbf{C}}(E,Y))$. \\

So let's rewrite \texttt{bind} in our notation:
\[
\begin{aligned}
	& (E\to X) \to & (X \to E \to Y) \to & E \to & Y \\
	& r \mapsto & f \mapsto & e \mapsto (f(r))(e)(e)
\end{aligned}
\]

To clarify these definitions, rewrite the reader monad: the endomorphism $T$ is defined as
\begin{equation}
\boxed{  \begin{aligned} & T := [E, -] \\ & T(X) = [E, X] \end{aligned} }
\end{equation}

For any morphism $f : X \to T(Y) = [E, Y]$, then \\
recall the Klesli composition: $g \circ_{\text{Klesli}} f = \mu_Z \circ Tg \circ f$ with $g: Y \to T(Z) = [E,Z]$. 

But if you take a look at how wikipedia defines "bind" for the reader monad, then bind \emph{is not} the Klesli composition, but rather it's $\mu_Y \circ Tf$ (!!!):
\begin{equation}
\begin{gathered}
\mu_Y \circ Tf: T(X) \to T(Y) \text{ or } \mu_Y \circ Tf : [E, X] \to [E, Y] \\ 
r \in [E,X] \qquad \, r(e) \in X \text{ for } e \in E \qquad \, f(r(e)) \in T(Y) = [E,Y] \qquad \, f(r(e))(e) \in Y
\end{gathered}
\end{equation}
bind, as defined by wikipedia for the reader monad \emph{is not} the Klesli composition.

As an example, consider the following example: \\
$T(Y) \ni (E\to Y) \in \text{Hom}_{\mathbf{C}}(E,Y)$ \\
$e = $ "dag" $\mapsto y = $ "Hello Dag!" \\

$T(X) \ni (E\to X) \in \text{Hom}_{\mathbf{C}}(E,X)$ \\
$e \equiv $ \texttt{name} $\mapsto x = $ "Hi \texttt{name}!" \\

$f:X\to T(Y) \to f:X\to \text{Hom}{(E,Y)}$ \\
$f(x) = $ replace "Hi" with "Hello" in $X$. \\

\texttt{ask}: $E\to E = 1_E$. \\
\texttt{ask} operation used to retrieve current context. \\

\texttt{local}: $(E\to E) \to (E\to T) \to E \to T = f \mapsto c \mapsto e \mapsto c(fe)$ \\
local executes a computation in a modified subcontext. $(E\to E)$ modifies the environment, $(E\to T)$ uses the modified environment. \\
Rewrite \texttt{local} in our notation: \\
\[
\begin{aligned} 
	& (E\to E) \to & (E\to X) \to & E \to & X \\
	& f \mapsto & c \mapsto & e \mapsto & c(f(e))
\end{aligned}
\]
To make sense of this, let $f \in [E, E]$, $c\in [E,X] = T(X)$, $c\circ f \in [E, X]$ so $(c\circ f)(e) \in X$. \\


\subsubsection{Writer Monad} $W \times (-) : \mathbf{C} \to \mathbf{C}$ where $W \in \text{Obj}{\mathbf{C}}$. If $W$ equipped with the structure of monoid, $W \times (-)$ canonically inherits the structure of a \emph{monad}. Hence, the Writer \emph{Monad}.  

Recall $W$ monoid if $W$ is a set equipped with binary operation $\mu : M \times M \to M$ s.t. $(xy)z = x(yz)$ and $1 \in M$ s.t. $1\cdot x = x = x \cdot 1$. \\

\subsubsection{State Monad}

cf. \href{https://ncatlab.org/nlab/show/state+monad}{state monad, nLab}

Input of type $X$, output of type $Y$, mutable state type $W$.

Consider morphism 
\[
X\times W \to Y\times W
\]

Under (Cartesian product $\vdash$ internal hom) adjunction, this is equivalently given by its adjunt, which is a function of type
\[
X \to [W, W\times Y]
\]
$[W, W\times (-)]$ is the monad, induced by the above adjunction. \\

$X \to [W, W\times Y]$ regarded as a morphism.

$[W, W\times Y]:\mathbf{H} \to \mathbf{H}$ is called \emph{state monad} for mutable states of type $W$.

TODO: Understand the above.

Writer comonad $W \times (-):\mathbf{C} \to \mathbf{C}$, \\

Given by forming Cartesian product with $W$. \\

Reader Monad $\circ$ Writer Comonad is state monad.

cf. \href{https://en.wikipedia.org/wiki/Monad_(functional_programming)#State_monads}{wikipedia}

unit or "return":
\begin{equation}\label{agat:Eq:ReturnStateMonadWikipedia}
\begin{aligned}
& X \to & S \to X \times S \\
& x \mapsto & (s \mapsto (x, s))
\end{aligned}
\end{equation} 
Compare this to 
\[
\begin{aligned} 
	& \eta_X : X \to T(X) = [S, S \times X] \\ 
	& \eta_X : x \mapsto (S \to S \times X)
\end{aligned} 
\]

"bind" was given by \href{https://en.wikipedia.org/wiki/Monad_(functional_programming)#State_monads}{wikipedia} as 
\begin{equation}\label{agat:Eq:BindStateMonadWikipedia}
\begin{aligned} 
	& (S \to X \times S)  \to & (X \to S \to Y \times S) \to & S \to & Y \times S \\ 
	& m \mapsto & k \mapsto & s \mapsto & k(x)(s) 
\end{aligned} 
\end{equation}
where $(x, t) = m(s)$

\subsubsection{Properties of the State Monad}

For the unit or "return", $\eta_X : X \to T(X) = [S, S \times X]$, \\

\begin{equation}
\begin{gathered} 
\eta_X : x \mapsto (S \to S \times X) \text{ s.t. } \\
(S \to S \times X) : s \mapsto (s, x)
\end{gathered} 
\end{equation}
Therefore, for the unit component at $X$ of a state monad,
\begin{equation}
\boxed{ \eta_X(x)(s) = (s, x) \text{ for } s \in S }
\end{equation}
Compare this to what wikipedia says how it defines the "\texttt{return}" operation in Eq. \ref{agat:Eq:ReturnStateMonadWikipedia}.

Consider any morphism $f : X \to T(Y) = [S, S\times Y]$,
\begin{equation}
\begin{aligned}
& f : X \to T(Y) = [S, S\times Y] \\
& f:x \mapsto (S \to S\times Y) \text{ s.t. } (S \to S \times Y) : s \mapsto (t = f_s(s, x), y = f_x(s, x))
\end{aligned}
\end{equation}

Allow morphism $f$ to have "freedom" to set state to arbitrary $t$.

\begin{equation}
\begin{aligned}
& Tf:T(X) \to T^2(Y) \to Tf : [S, S\times X] \to T([S, S\times Y]) = [S, S \times [S, S\times Y]] \\
& Tf:(S \to S \times X) \mapsto (S \to S \times (S \to S \times Y))
\end{aligned}
\end{equation}

Then
\begin{equation}
\begin{aligned}
& \mu_Y \circ Tf : T(X) \to T(Y) \text{ or } \mu_Y \circ Tf : [S, S \times X] \to [S, S \times Y] \\ 
& \mu_Y \circ Tf: (S \to S \times X) \mapsto (S \to S \times Y)
\end{aligned}
\end{equation}

Note that 
\begin{equation}
	\begin{aligned}
	& \mu_X : T^2(X) \to T(X) \\ 
	& \mu_X : (S \to S \times (S \to S \times X)) \mapsto (S \to S \times X)
	\end{aligned}
\end{equation} 

\begin{equation}
\begin{aligned}
& T^2 f : T^2(X) \to T^3(Y) \\  
& T^2 f: (S \to S \times (S \to S\times X)) \mapsto (S \to S \times (S \to S \times (S \to S \times Y)))
\end{aligned}
\end{equation}

Now
\[
\begin{gathered} 
 \mu_X \circ T^2 f : (S \to (S \times (S \to S \times X))) \mapsto (S \to S \times (S \to S \times Y)) \\ 
 Tf \circ \mu_X : (S \to S \times (S \to S \times X)) \mapsto (S \to S \times (S \to S \times Y))
\end{gathered} 
\]
and so 
\[
\Longrightarrow \mu_X T^2 = Tf \mu_X 
\]
Thus, $\mu$, as defined above, \emph{is} a natural transformation.

Show $\eta_{T(Y)} \circ f = Tf\circ \eta_X$. 
\[
\begin{gathered} 
	Tf \circ \eta_X (x) = (S \to S \times (S \to S \times Y)) \\ 
 \eta_{T(Y)} \circ f(x) = \eta_{T(Y)} (S \to S \times Y) = (S \to S \times (S \to S \times Y))
\end{gathered} 
\]


Prove now the coherence conditions, Prop. \ref{agat:Prop:MonadCoherenceConditions}. For the first one,
\[
\mu \circ T \mu = \mu \circ \mu T
\]

\[
\begin{gathered} 
	\mu \circ T \mu (S \to (S \to S \times (S\to S \times X))) = \mu \circ T (S \times (S \to S \times X)) = \mu (S \to S \times (S \times (S \to S \times Y))) = \\
	= (S \times (S \to S \times Y)) \\ 
 \mu\circ \mu T (S \to (S \to S \times (S\to S \times X))) = \mu^2 (S \to S \times (S \to (S \to S \times (S \to S \times Y)))) = \\
 	= S \times (S \to S \times Y)
\end{gathered} 
\]
\[
\Longrightarrow \mu \circ T \mu = \mu \circ \mu T
\]
So the first coherence condition is proven.

For the second, 
\[
\begin{gathered}
\mu \circ T \eta (S\to S \times X) = \mu \circ T (S \to S \times (S \to S \times X)) = (S \to S \times (S \to S \times X)) \\
\mu \eta T (S\to S \times X) = \mu \circ (S \to S \times T(S \to S \times X)) = T(S \to S \times X) = (S \to S \times (S \to S \times X))
\end{gathered}
\]

\[
\]
Thus, it's been shown that the choice of endomorphism for the state monad, $X \to [S, S \times X]$, fulfills the monad coherence conditions being of this form.












Consider what "bind" means in category theory. Given
\begin{equation}
\begin{aligned}
& f:X \to T(Y) = [S, S \times Y] \\ 
& g: Y \to T(Z) = [S, S\times Z] \\ 
& f(x) = (S \to S\times Y) \\
& f(x)(s) = (t, y) \text{ where } s, t \in W \\ 
& g(y) = (S \to S \times Z) \\ 
& g(y)(t) = (u, z)
\end{aligned}
\end{equation}
Therefore, if we define the usual projection operations on a Cartesian product,
\begin{equation}
\begin{aligned}
& \Pi_W : W \times X \to W \\
& \Pi_W(w,x) = w
\end{aligned} \qquad \, 
\begin{aligned}
& \Pi_X : W \times X \to X \\
& \Pi_X(w,x) = x
\end{aligned}
\end{equation}
\begin{equation}
\begin{aligned}
& \mu_Z \circ Tg \circ f : X \to T(Z) = [S, S\times Z] \ni (W \to W \times Z) \\ 
& \mu_Z \circ Tg \circ f : x \mapsto (S \to S \times Z) \equiv \mu_Z \circ Tg \circ f(x)
\end{aligned}
\end{equation}
such that 
\begin{equation}
\boxed{ g(y)(t) = (u,z) \text{ where } y = \Pi_Y(f(x)(s)), \, t = \Pi_W(f(x)(s)) }
\end{equation}
Compare this to what wikipedia says in Eq. \ref{agat:Eq:BindStateMonadWikipedia}.

\subsubsection{Continuation Monad}

From \href{https://ncatlab.org/nlab/show/continuation+monad}{nLab}, \\
In a category $\mathbf{C}$ with internal homs $[-, -]$, so internal homs $[-,-]$ belong to $\text{Obj}{\mathbf{C}}$ (they're objects as well), given object $Y \in \text{Obj}{\mathbf{C}}$, the \emph{continuation monad} is endofunctor
\begin{equation}
\boxed{ X \mapsto [[X, Y], Y] } \text{ i.e. } T(X) = [[X, Y], Y] \text{ for object $Y \in \text{Obj}{\mathbf{C}}$ } 
\end{equation}
so that $T = X\mapsto [[X, Y] , Y]$

From \href{https://en.wikipedia.org/wiki/Monad\_(functional\_programming)\#Continuation\_monad}{wikipedia},

unit or return:  \\

return type $Y$, maps type $X$ into functions of type $(X \to Y) \to Y$
\begin{equation} 
\begin{aligned}
	& X \to & (X\to Y) \to & Y \\
	& x\mapsto & f \mapsto & f(x) = y
\end{aligned}
\end{equation}

Compare this to
\begin{equation}
\begin{aligned}
& T : X \to [[X, Y], Y] \\ 
& \eta_X : X \to T(X) = [[X,Y], Y]
\end{aligned}
\end{equation}

And so for the \emph{unit} of the Continuation monad,
\begin{equation}
\boxed{
\begin{aligned}
& \eta_X : X \to T(X) = [[X, Y], Y] \\
& \eta_X : x\mapsto (f \to y) \\ 
& \eta_X(x)(f) = y = f(x)
\end{aligned}
}
\end{equation}

bind: \\

\begin{equation}\label{agat:Eq:ContinuationMonadBindWiki}
\begin{aligned}
& ((X\to Y) \to Y) \to & (X \to (X' \to Y) \to Y) \to & (X' \to Y) \to & Y \\
& c \mapsto & f \mapsto & k \mapsto &  c(t\mapsto ftk)
\end{aligned}
\end{equation}

From \href{https://ncatlab.org/nlab/show/continuation+monad}{nLab}, if morphism $f:X\to Y$ is in Klesi category of the continuation monad, then it's a morphism in $\mathbf{C}$ of form $X\to T(Y)$ in $\mathbf{C}$. Hence morphism in original category \emph{of the form} $X\to [[Y,S], S]$ is much like a map from $X \to Y$, only instead of "returning" output $Y$ directly, it instead feeds it into given function $Y\to S$, which hence \textbf{continues} the computation. 

So for $X \to [[Y, S], S]$, \href{https://ncatlab.org/nlab/show/internal+hom}{nLab} gives the corresponding \emph{computation rule} for function types of the internal hom $[X, S]$:
\[
(y \mapsto a(y))(x) = a(x)
\]

Hence for $X \to [[X, Y], Y]$, we have $x\mapsto f \mapsto f(x)$ because $f \in [X, Y]$ because $f$ is treated as an object in $\mathbf{C}$. \\

So do the following "rewrite" or change of notation on Eq. \ref{agat:Eq:ContinuationMonadBindWiki}:

\begin{equation}\label{agat:Eq:ContinuationMonadBindWiki2}
\begin{aligned}
& ((X \to Z) \to Z) \to (X \to (Y\to Z) \to Z) \to (Y \to Z) \to Z \\ 
& T(X) \to (X \to T(Y)) \equiv \text{Hom}(X, T(Y)) \to \text{Hom}(Y,Z) \to Z
\end{aligned}
\end{equation}
Let $c\in T(X) = [[X,Z], Z]$. \\
Let $k \in [Y,Z]$ \\
Let $f(x) \in [[Y,Z], Z]$ and so \\
$f(x)(k) \in Z$ \\

Form $x \mapsto f(x)(k) \in [X,Z]$, so $c(x\mapsto f(x)(k)) \in Z$. \\

So "bind" as defined in wikipedia, Haskell/Functional Programming isn't the Klesli composition, but it's
\[
\mu_Y \circ Tf : T(X) \to T(Y) = [[Y,Z], Z]
\]

Write this as 
\begin{equation}
\boxed{
\begin{aligned}
& (\mu_Y \circ T)(f)(c) \in T(Y) = [[Y,Z], Z] \\
& (\mu_Y \circ T)(f)(c)(k) \in Z
\end{aligned}}
\end{equation}


\emph{call with current continuation} function defined as follows: \\
call/cc
\begin{equation}\label{agat:Eq:CallCurrentContinuationContinuationMonad}
\begin{aligned}
((X \to (X' \to Y) \to Y) \to (X\to Y) \to Y) & \to (X\to Y)  & \to Y \\ 
f & \mapsto k & \mapsto (f(t\mapsto x \mapsto kt) k) 
\end{aligned}
\end{equation}

Let's look at Eq. \ref{agat:Eq:CallCurrentContinuationContinuationMonad}, step-by-step.
\[
\begin{gathered}
\begin{aligned} 
& (X\to (X' \to Y) \to Y) \\
& f
\end{aligned} \text{ corresponds to } \\
\begin{gathered}
X \to [[X' , Y], Y] = T(X') \\
f \equiv f_{XX'} : X \to T(X')
\end{gathered}
\end{gathered}
\]

\subsubsection{\texttt{future} as Continuation Monad}

 \v{C}uki\'{c} (2018) \cite{agat:Cuki2018} asks first if \verb|std::future| in C++ is a functor. Recall the definition of a functor, Def. \ref{agat:Def:Functor}:

\[
\begin{aligned}
& F: \mathbf{C} \to \mathbf{D} \text{ with "object" map } \\ 
& F: \text{Obj}{\mathbf{C}} \to \text{Obj}{\mathbf{D}} \\ 
& F: X \to F(X)
\end{aligned}
\] 
and the functor's "morphism map":
\[
\begin{aligned}
& Ff :F(X) \to F(Y) \\ 
& f:X \to Y, \, f \in \text{Hom}_{\mathbf{C}}(X, Y)
\end{aligned}
\]

Suppose $X = $ \texttt{T1}, $Y = $ \texttt{T2}. Then in the example on pp. 220, Sec. 10.7.1 Futures as monads" example of \v{C}uki\'{c} (2018) \cite{agat:Cuki2018}, 
\[
\begin{aligned}
& F(X) = \verb|future<T1>| \\ 
& F(Y) = \verb|future<T2>|
\end{aligned}
\]

Instead, the morphisms to consider is of this form:
\[
f:X \to T(Y)
\]
where $T$ is a functor but a special one: the endomorphism. And $T$ is \verb|std::future|.

Here are examples of morphisms $f, g$:
\[
\begin{aligned}
& \verb|future<std::string> user_full_name(const std::string& login);| \\
& \verb|future<std::string> to_html(const std::string& text);|
\end{aligned}
\]

Now $Tf:T(X) \to T^2(Y)$, how the morphism map for $T$ applies to morphism $f$, now. Recall that we can compose this with the multiplication natural transformation component at $X$, $\mu_X$:
\[
\begin{aligned}
& \mu_X \circ Tf : T(X) \to T(Y) \\ 
& \mu_X : T^2(X) \to T(X)
\end{aligned}
\]

In order to compose morphisms, say $f,g$, recall the "Klesli composition", Prop. \ref{agat:Prop:KlesliCompositionForm}
\begin{equation}
g\circ_{\text{Klesili}} f = \mu_Z \circ Tg \circ f
\end{equation}
If we call $\mu_Z \circ T$ "bind", $\mu_Z \circ T: \text{Hom}(Y, T(Z)) \to \text{Hom}(T(Y), T(Z))$. The correspondence is such:
\begin{equation}
(\texttt{bind}(\verb|to_html|))(\verb|user_full_name|) \Longleftrightarrow \mu_Z \circ Tg \circ f: X \to T(Z)
\end{equation}


\subsubsection{Resources for implementations of Monads} Python and C++ code. \\

\url{https://github.com/Iasi-C-CPP-Developers-Meetup/presentations-code-samples}, namely \\
\url{https://github.com/Iasi-C-CPP-Developers-Meetup/presentations-code-samples/tree/master/radugeorge} \\

\url{https://github.com/dbrattli/OSlash}, namely \\
\url{https://github.com/dbrattli/OSlash/tree/master/tests} \\

\href{https://gist.github.com/pfultz2/f538423ebc9bbbe3df92}{State Monad implementation in C++}


\section{Applications of Category Theory on Hybrid Systems}

cf. Ames (2006) \cite{agat:Ames2006}.
	
\subsection{D-Categories}

$D$ stands for discrete.

Recall that a small category $\mathbf{C}$ is called \emph{small} if both $\text{Obj}(\mathbf{C})$ and $\text{hom}(\mathbf{C})$ are sets, not proper classes. 

\begin{definition}[Axiomatic D-categories]
Let D-category be a \emph{small} category $\mathbf{D}$ s.t.
\begin{enumerate}
	\item $\forall \, D \in \text{Obj}(\mathbf{D})$, \\
	$\exists \,$ morphism $f \in \text{Mor}(\mathbf{D})$ s.t. $f \neq 1$ s.t. \\
	$f \in \text{Hom}(D, *)$ or $f\in \text{Hom}(*, D)$, but never both, \\
	i.e. $\forall \, $ diagram $a_0 \xrightarrow{ \alpha_1} a_1 \xrightarrow{ \alpha_2} \dots \xrightarrow{ \alpha_n} a_n$ in $\mathbf{D}$, all but 1 morphism must be identity (i.e. longest chain of composite non-identity morphisms is of length 1). 
	\item If for $D\in \text{Obj}(\mathbf{D})$, $D$ is the domain of a non-identity morphism, i.e. $\exists \, f_1 \in \text{Mor}(\mathbf{D})$ s.t. 
	\[
	f_1 \in \text{Hom}(D, *), \, f_1 \neq 1
	\]
	then $\exists \, f_2 \in \text{Hom}(D, *)$, $f_2 \neq f_1$, $f_2 \neq 1$ and $\forall \, f \in \text{Hom}(D, *)$ s.t. $f\neq f_1,f_2$, \, $f=1$
\end{enumerate}
\end{definition}
	
cf. 1.2.1 Important objects in D-categories, Ames (2006) \cite{agat:Ames2006}.

Let 
\[
\begin{gathered}
	\text{Mor}{\mathbf{A}} = \bigcup_{A, B \in \text{Obj}(\mathbf{A})} \text{Hom}(A,B) \qquad \, \text{(my notation)} \\
	\text{Mor}(\mathcal{D}) = \bigcup_{(a,b) \in \text{Obj}(\mathcal{D}) \times \text{Obj}(\mathcal{D}) } \text{Hom}_{\mathcal{D}}(a,b) \qquad \, \text{ (Ames' notation) }
\end{gathered}
\]	
	
Let 
\begin{equation}
\text{Mor}_{1}{\mathbf{A}} := \lbrace A \in \text{Mor}(\mathbf{A}) | A \neq 1 \rbrace 
\end{equation}

For a D-category, consider these subset of $\text{Obj}(\mathbf{D})$, 

\begin{definition}[Edge set]
	\textbf{edge set} of $\mathbf{D}$, $E(\mathbf{D})$, 
\begin{equation}
\begin{gathered}
E(\mathbf{D}) := \lbrace A \in \text{Obj}(\mathbf{D}) | \alpha \in \text{Hom}(A, *), \beta \in \text{Hom}(A, *), \alpha, \beta \in \text{Mor}_{1}(\mathbf{D}), \, \alpha \neq \beta \rbrace \text{ i.e. } \\
E(\mathbf{D}) := \lbrace A \in \text{Obj}(\mathbf{D}) | \alpha, \beta \in \text{Hom}(A, *), \, \alpha, \beta \neq 1, \, \alpha \neq \beta \rbrace
\end{gathered}
\end{equation}	
\end{definition}
i.e. $\forall \, A \in E(\mathbf{D})$, $\exists \, \alpha, \beta \in \text{Mor}(\mathbf{D})$, \, $\alpha, \beta \neq 1$, s.t. $\alpha, \beta \in \text{Hom}(D, *)$; denote these morphisms by $s_a,t_a$ (this specific choice will define an \textbf{orientation}).

Conversely, given morphism $\gamma \in \text{Mor}(\mathbf{D})$, $\gamma \neq 1$, $\exists \, ! \, A \in E(\mathbf{D})$ s.t. $\gamma = s_a$ or $\gamma = t_a$, i.e. $\gamma \in \text{Hom}(A, *)$. 

\begin{definition}[Vertex set]
\textbf{vertex set} of $\mathbf{D}$:
\begin{equation}
 V(\mathbf{D}) = (E(\mathbf{D}))^c
\end{equation} 
\end{definition}

\begin{definition}[Orientation]
	Orientation of D-category $\mathbf{D}$ is a pair of functions $(s,t)$ between sets.
	
\end{definition}


	


